% app-glossary-generated.tex -- GENERATED by docs/gen-glossary.py from
% reference/GLOSSARY.md.  Do not edit; re-run `make gen' in docs/.

\chapter{Glossary: every name, A to Z}
\label{app:glossary}

Every other listing in this manual and in the browser reference is grouped by
something: operators by class, structures by structure, tactics by task.  This
one has the ordering nobody has to learn.  It carries the whole
\emph{vocabulary} -- 516 names: 16 structure, 7 refinement, 8 instance, 17 view, 163 predicate, 112 functoid, 25 accessor, 18 defined-fn, 60 operator, 8 primitive, 82 tactic -- and no results: the theorems, axioms and
supports are catalogued in \texttt{THEOREMS.md}, \texttt{BY-OPERATOR.md} and
\texttt{PROOF-DEBT.md}, and each entry here says how many of them mention it.

At the REPL, \texttt{(glossary 'NAME)} prints one entry and
\texttt{(glossary "pat")} every name containing \texttt{pat}.

\begingroup
\small
\setlength{\parindent}{0pt}

\section*{Symbols}
\addcontentsline{toc}{section}{Symbols}

\noindent\textbf{\texttt{*}}\quad\textit{operator}\\
a kernel term-former \\
mentioned by 238 result(s)

\smallskip

\noindent\textbf{\texttt{+}}\quad\textit{operator}\\
a kernel term-former \\
mentioned by 190 result(s)

\smallskip

\noindent\textbf{\texttt{-}}\quad\textit{operator}\\
a kernel term-former \\
mentioned by 131 result(s)

\smallskip

\noindent\textbf{\texttt{/}}\quad\textit{operator}\\
a kernel term-former \\
mentioned by 16 result(s)

\smallskip

\noindent\textbf{\texttt{<(x, y)}}\quad\textit{predicate}\\
reads: \$1 is less than \$2 \\
defined by \texttt{<} \\
mentioned by 129 result(s)

\smallskip

\noindent\textbf{\texttt{<=}}\quad\textit{primitive}\\
reads: \$1 is at most \$2 \\
a kernel relation \\
mentioned by 300 result(s)

\smallskip

\noindent\textbf{\texttt{<=\_ord}}\quad\textit{predicate}\\
a predicate \\
mentioned by 16 result(s)

\smallskip

\noindent\textbf{\texttt{<\_ord}}\quad\textit{predicate}\\
a predicate \\
mentioned by 13 result(s)

\smallskip

\noindent\textbf{\texttt{=}}\quad\textit{primitive}\\
reads: \$1 equals \$2 \\
a kernel relation

\smallskip

\noindent\textbf{\texttt{==}}\quad\textit{primitive}\\
reads: \$1 is identical to \$2 \\
a kernel relation

\smallskip

\noindent\textbf{\texttt{>}}\quad\textit{primitive}\\
reads: \$1 is greater than \$2 \\
a kernel relation

\smallskip

\noindent\textbf{\texttt{>=}}\quad\textit{primitive}\\
reads: \$1 is at least \$2 \\
a kernel relation

\smallskip


\section*{A}
\addcontentsline{toc}{section}{A}

\noindent\textbf{\texttt{abelian-group}}\quad\textit{structure}\\
predicate \texttt{is-abelian-group}; slots: carr opr iden inv

\smallskip

\noindent\textbf{\texttt{abelian-group-as-monoid(s)}}\quad\textit{view}\\
an abelian-group viewed as a monoid \\
mentioned by 49 result(s)

\smallskip

\noindent\textbf{\texttt{abs}}\quad\textit{operator}\\
a kernel term-former \\
mentioned by 60 result(s)

\smallskip

\noindent\textbf{\texttt{act(s)}}\quad\textit{accessor}\\
a structure slot \\
mentioned by 95 result(s)

\smallskip

\noindent\textbf{\texttt{add(s)}}\quad\textit{accessor}\\
a structure slot \\
mentioned by 420 result(s)

\smallskip

\noindent\textbf{\texttt{ai}}\quad\textit{tactic}\\
\texttt{(ai hyp)} -- Antecedent inference: decompose a cited assumption -- AND-split, OR-into-cases, or FORSOME-elimination to a fresh eigenvariable. \\
declared in \texttt{tactics-help.scm}

\smallskip

\noindent\textbf{\texttt{apply-functoid}}\quad\textit{operator}\\
a kernel term-former

\smallskip

\noindent\textbf{\texttt{arith}}\quad\textit{tactic}\\
\texttt{(arith)} -- Discharge a ground arithmetic goal by evaluation. \\
declared in \texttt{tactics-help.scm}

\smallskip

\noindent\textbf{\texttt{ass}}\quad\textit{tactic}\\
\texttt{(ass)} -- Close the goal by an assumption alpha-equivalent to it. \\
declared in \texttt{tactics-help.scm}

\smallskip

\noindent\textbf{\texttt{ass-all-frontier!}}\quad\textit{tactic}\\
\texttt{(ass-all-frontier!)} -- Close every frontier leaf whose goal is already among its assumptions.  [proof-local] \\
declared in \texttt{tactics-help.scm}

\smallskip


\section*{B}
\addcontentsline{toc}{section}{B}

\noindent\textbf{\texttt{ball(s, c, r)}}\quad\textit{functoid}\\
= \{y in pts(s): (dist(s))(c, y) <= r and not((dist(s))(c, y) = r)\} \\
mentioned by 20 result(s)

\smallskip

\noindent\textbf{\texttt{ball-cover(s, r)}}\quad\textit{functoid}\\
= image(vnb-lambda(c, pts(s), ball(s, c, r)), pts(s)) \\
mentioned by 4 result(s)

\smallskip

\noindent\textbf{\texttt{bc}}\quad\textit{tactic}\\
\texttt{(bc impl)} -- Backchain the goal through an (IMPLIES A B) already in context: if the goal matches B, the new goal is A. \\
declared in \texttt{tactics-help.scm}

\smallskip

\noindent\textbf{\texttt{bc*}}\quad\textit{tactic}\\
\texttt{(bc* 'thm [((v val)...)] h1 h2 ...)} -- Matching backchain: unify the theorem's conclusion against the goal, then leave its (instantiated) antecedents as subgoals -- optional handlers hk run on the k-th subgoal. \\
declared in \texttt{tactics-help.scm}

\smallskip

\noindent\textbf{\texttt{bdd-metric(s)}}\quad\textit{functoid}\\
= [pts(s), vnb-lambda([u, v], cartesian(pts(s), pts(s)), /((dist(s))(u, v), 1 + (dist(s))(u, v)))] \\
mentioned by 11 result(s)

\smallskip

\noindent\textbf{\texttt{beta}}\quad\textit{tactic}\\
\texttt{(beta)} -- Beta-reduce a functoid application in the goal. \\
declared in \texttt{tactics-help.scm}

\smallskip

\noindent\textbf{\texttt{big-union}}\quad\textit{operator}\\
a kernel term-former \\
mentioned by 14 result(s)

\smallskip

\noindent\textbf{\texttt{bijection}}\quad\textit{operator}\\
a kernel term-former \\
mentioned by 36 result(s)

\smallskip

\noindent\textbf{\texttt{binneg}}\quad\textit{operator}\\
a kernel term-former \\
mentioned by 6 result(s)

\smallskip

\noindent\textbf{\texttt{binplus}}\quad\textit{operator}\\
a kernel term-former \\
mentioned by 14 result(s)

\smallskip

\noindent\textbf{\texttt{bintimes}}\quad\textit{operator}\\
a kernel term-former \\
mentioned by 10 result(s)

\smallskip

\noindent\textbf{\texttt{block(p, k, l)}}\quad\textit{functoid}\\
= matof(k, l, vnb-lambda([i, j], cartesian(interval(1, k), interval(1, l)), entry(p, i, j))) \\
mentioned by 7 result(s)

\smallskip

\noindent\textbf{\texttt{border(a, b, m, p, q)}}\quad\textit{functoid}\\
= matof(succ(p), succ(q), vnb-lambda([i, j], cartesian(interval(1, succ(p)), interval(1, succ(q))), if(i = 1, if ... \\
mentioned by 22 result(s)

\smallskip

\noindent\textbf{\texttt{bu-me}}\quad\textit{tactic}\\
\texttt{(bu-me hyp)} -- Big-union membership elim. \\
declared in \texttt{tactics-help.scm}

\smallskip

\noindent\textbf{\texttt{bu-mi}}\quad\textit{tactic}\\
\texttt{(bu-mi w)} -- Big-union membership intro via the index witness w. \\
declared in \texttt{tactics-help.scm}

\smallskip

\noindent\textbf{\texttt{bu-set}}\quad\textit{tactic}\\
\texttt{(bu-set)} -- Big-union sethood. \\
declared in \texttt{tactics-help.scm}

\smallskip


\section*{C}
\addcontentsline{toc}{section}{C}

\noindent\textbf{\texttt{calc}}\quad\textit{tactic}\\
\texttt{(calc L0 (rel1 L1 [just1]) (rel2 L2 [just2]) ...)} -- Ground the focus goal (REL L0 Ln) by a CHAIN of intermediaries L0 rel1 L1 rel2 L2 ... reln Ln: proves each link and composes them into the endpoint.  A link no lane can close is LEFT OPEN as a leaf -- the refinement point, and the PSS candidate.  Relations: = == (folded by transitivity), < <= (folded through the co-*-trans lemmas, with = steps riding along), IFF (chained per direction).  RETURNS the links, the open ones, and their PSS candidates as an alist. \\
declared in \texttt{tactics-help.scm}

\smallskip

\noindent\textbf{\texttt{card}}\quad\textit{operator}\\
a kernel term-former \\
mentioned by 153 result(s)

\smallskip

\noindent\textbf{\texttt{carr(s)}}\quad\textit{accessor}\\
a structure slot \\
mentioned by 1072 result(s)

\smallskip

\noindent\textbf{\texttt{cartesian}}\quad\textit{operator}\\
a kernel term-former \\
mentioned by 95 result(s)

\smallskip

\noindent\textbf{\texttt{cauchy-setoid(m)}}\quad\textit{functoid}\\
= [cseq(m), crel(m)] \\
mentioned by 1 result(s)

\smallskip

\noindent\textbf{\texttt{cc-ms}}\quad\textit{instance}\\
an instance of metric-space; predicate \texttt{is-cc-ms}; same shape as metric-space

\smallskip

\noindent\textbf{\texttt{cc-normed-field}}\quad\textit{instance}\\
an instance of normed-field; predicate \texttt{is-cc-normed-field}; same shape as normed-field

\smallskip

\noindent\textbf{\texttt{ccint(a, b)}}\quad\textit{functoid}\\
= \{x in rr: a <= x and x <= b\} \\
mentioned by 27 result(s)

\smallskip

\noindent\textbf{\texttt{ce}}\quad\textit{tactic}\\
\texttt{(ce hyp k)} -- Cartesian elim: project the k-th component out of a CARTESIAN-membership assumption. \\
declared in \texttt{tactics-help.scm}

\smallskip

\noindent\textbf{\texttt{centre-set(s, r, f)}}\quad\textit{functoid}\\
= image(vnb-lambda(b, f, choice(centres(s, b, r))), f) \\
mentioned by 2 result(s)

\smallskip

\noindent\textbf{\texttt{centres(s, b, r)}}\quad\textit{functoid}\\
= \{c in pts(s): ball(s, c, r) = b\} \\
mentioned by 6 result(s)

\smallskip

\noindent\textbf{\texttt{choice}}\quad\textit{operator}\\
a kernel term-former \\
mentioned by 9 result(s)

\smallskip

\noindent\textbf{\texttt{choose(n, m)}}\quad\textit{functoid}\\
= card(choose-set(n, m)) \\
mentioned by 9 result(s)

\smallskip

\noindent\textbf{\texttt{choose-set(n, m)}}\quad\textit{functoid}\\
= \{a in power(ord-segment(n)): card(a) = m\}

\smallskip

\noindent\textbf{\texttt{ci}}\quad\textit{tactic}\\
\texttt{(ci)} -- Cartesian intro: prove a CARTESIAN-product membership component-wise. \\
declared in \texttt{tactics-help.scm}

\smallskip

\noindent\textbf{\texttt{class(s, a)}}\quad\textit{functoid}\\
= \{b in pts(s): related(s, a, b)\} \\
mentioned by 25 result(s)

\smallskip

\noindent\textbf{\texttt{closure(s, a)}}\quad\textit{functoid}\\
= \{x in pts(s): forall([u], is-open(s, u) and x in u implies forsome([y in u], y in a))\} \\
mentioned by 2 result(s)

\smallskip

\noindent\textbf{\texttt{cluster-point(s, f, x)}}\quad\textit{predicate}\\
reads: \$3 is a cluster point of the sequence \$2 in \$1 \\
defined by \texttt{cluster-point} \\
mentioned by 6 result(s)

\smallskip

\noindent\textbf{\texttt{comb-kk}}\quad\textit{defined-fn}\\
defined by \texttt{comb-kk-zero} \texttt{comb-kk-succ} \\
mentioned by 13 result(s)

\smallskip

\noindent\textbf{\texttt{comm-monoid}}\quad\textit{structure}\\
predicate \texttt{is-comm-monoid}; slots: carr opr iden

\smallskip

\noindent\textbf{\texttt{commutative-ring}}\quad\textit{refinement}\\
predicate \texttt{is-commutative-ring}; same shape as ring

\smallskip

\noindent\textbf{\texttt{commutative-ring-additive-ag(s)}}\quad\textit{view}\\
a commutative-ring viewed as an abelian-group \\
mentioned by 20 result(s)

\smallskip

\noindent\textbf{\texttt{commutative-ring-multiplicative-cm(s)}}\quad\textit{view}\\
a commutative-ring viewed as a comm-monoid \\
mentioned by 5 result(s)

\smallskip

\noindent\textbf{\texttt{comp}}\quad\textit{operator}\\
a kernel term-former

\smallskip

\noindent\textbf{\texttt{comp-me}}\quad\textit{tactic}\\
\texttt{(comp-me hyp)} -- Comprehension membership elim. \\
declared in \texttt{tactics-help.scm}

\smallskip

\noindent\textbf{\texttt{comp-mi}}\quad\textit{tactic}\\
\texttt{(comp-mi)} -- Comprehension membership intro. \\
declared in \texttt{tactics-help.scm}

\smallskip

\noindent\textbf{\texttt{complement-in}}\quad\textit{operator}\\
a kernel term-former \\
mentioned by 11 result(s)

\smallskip

\noindent\textbf{\texttt{completion(m)}}\quad\textit{functoid}\\
= [quotient(cauchy-setoid(m)), completion-dist(m)] \\
mentioned by 5 result(s)

\smallskip

\noindent\textbf{\texttt{completion-dist(m)}}\quad\textit{functoid}\\
= vnb-lambda(p, cartesian(quotient(cauchy-setoid(m)), quotient(cauchy-setoid(m))), iota(dval, forsome([f, g], f  ...

\smallskip

\noindent\textbf{\texttt{compose(f, g)}}\quad\textit{functoid}\\
= vnb-lambda(z\_, dom(g), f(g(z\_))) \\
mentioned by 11 result(s)

\smallskip

\noindent\textbf{\texttt{conjugate}}\quad\textit{operator}\\
a kernel term-former \\
mentioned by 13 result(s)

\smallskip

\noindent\textbf{\texttt{converges(s, f)}}\quad\textit{predicate}\\
reads: \$2 converges in \$1 \\
defined by \texttt{converges} \\
mentioned by 14 result(s)

\smallskip

\noindent\textbf{\texttt{converges-along(s, g, b, p)}}\quad\textit{predicate}\\
reads: \$2 converges to \$4 along \$3 in \$1 \\
defined by \texttt{converges-along} \\
mentioned by 4 result(s)

\smallskip

\noindent\textbf{\texttt{converges-on(s, fam, dseq)}}\quad\textit{predicate}\\
reads: \$2 converges at every point of the sequence \$3 in \$1 \\
defined by \texttt{converges-on} \\
mentioned by 4 result(s)

\smallskip

\noindent\textbf{\texttt{converges-to(s, f, l)}}\quad\textit{predicate}\\
reads: \$2 converges to \$3 in \$1 \\
defined by \texttt{converges-to} \\
mentioned by 19 result(s)

\smallskip

\noindent\textbf{\texttt{converges-uniformly(s, seq, g)}}\quad\textit{predicate}\\
reads: \$2 converges uniformly to \$3 on \$1 \\
defined by \texttt{converges-uniformly} \\
mentioned by 5 result(s)

\smallskip

\noindent\textbf{\texttt{cos}}\quad\textit{operator}\\
a kernel term-former

\smallskip

\noindent\textbf{\texttt{crel(m)}}\quad\textit{functoid}\\
= \{p in cartesian(cseq(m), cseq(m)): cseq-equiv(m, nth(1, p), nth(2, p))\}

\smallskip

\noindent\textbf{\texttt{crs}}\quad\textit{tactic}\\
\texttt{(crs)} -- Commutative-ring decision procedure: prove a polynomial identity over ZZ[generators].  Expands literal powers, so (x+y)\textasciicircum{}2 = ... closes directly. \\
declared in \texttt{tactics-help.scm}

\smallskip

\noindent\textbf{\texttt{cseq(m)}}\quad\textit{functoid}\\
= \{f in fun(nn, pts(m)): is-cauchy-seq(m, f)\}

\smallskip

\noindent\textbf{\texttt{cseq-equiv(m, f, g)}}\quad\textit{predicate}\\
reads: \$2 and \$3 are equivalent Cauchy sequences in \$1 \\
defined by \texttt{cseq-equiv} \\
mentioned by 2 result(s)

\smallskip

\noindent\textbf{\texttt{cut}}\quad\textit{tactic}\\
\texttt{(cut formula)} -- Cut: prove \texttt{formula' as a side subgoal, then continue the main goal with }formula' added as an assumption (Gentzen cut). \\
declared in \texttt{tactics-help.scm}

\smallskip

\noindent\textbf{\texttt{cut-mem!}}\quad\textit{tactic}\\
\texttt{(cut-mem! mem A)} -- Prove a membership (IN (f x) B) by fun-apply-type with domain A, leaving it in context.  [proof-local] \\
declared in \texttt{tactics-help.scm}

\smallskip


\section*{D}
\addcontentsline{toc}{section}{D}

\noindent\textbf{\texttt{delete-at}}\quad\textit{operator}\\
a kernel term-former \\
mentioned by 6 result(s)

\smallskip

\noindent\textbf{\texttt{deriv(f, a)}}\quad\textit{functoid}\\
= iota(l, is-diff-at(f, a, l)) \\
mentioned by 4 result(s)

\smallskip

\noindent\textbf{\texttt{deriv-v(m, f, a)}}\quad\textit{functoid}\\
= iota(l, is-diff-at-v(m, f, a, l)) \\
mentioned by 2 result(s)

\smallskip

\noindent\textbf{\texttt{descend(s, f)}}\quad\textit{functoid}\\
= vnb-lambda(c, quotient(s), iota(z, forsome([a in c], z = f(a)))) \\
mentioned by 5 result(s)

\smallskip

\noindent\textbf{\texttt{descend2(s, f)}}\quad\textit{functoid}\\
= vnb-lambda([c, d], cartesian(quotient(s), quotient(s)), iota(z, forsome([a in c, b in d], z = f(a, b)))) \\
mentioned by 7 result(s)

\smallskip

\noindent\textbf{\texttt{det}}\quad\textit{primitive}\\
reads: det(\$3) \\
defined by \texttt{det-zero} \texttt{det-cofactor} \\
mentioned by 15 result(s)

\smallskip

\noindent\textbf{\texttt{detach!}}\quad\textit{tactic}\\
\texttt{(detach! impl)} -- Forward modus ponens: from an in-context (IMPLIES A B) whose A is also in context, leave B in context. \\
declared in \texttt{tactics-help.scm}

\smallskip

\noindent\textbf{\texttt{di}}\quad\textit{tactic}\\
\texttt{(di)} -- Direct inference: split an AND goal, move an IMPLIES antecedent into the assumptions, or introduce a fresh eigenvariable for a leading FORALL. \\
declared in \texttt{tactics-help.scm}

\smallskip

\noindent\textbf{\texttt{difference}}\quad\textit{operator}\\
a kernel term-former \\
mentioned by 13 result(s)

\smallskip

\noindent\textbf{\texttt{dist(s)}}\quad\textit{accessor}\\
a structure slot \\
mentioned by 76 result(s)

\smallskip

\noindent\textbf{\texttt{dist-seq(m, f, g)}}\quad\textit{functoid}\\
= vnb-lambda(n, nn, (dist(m))(f(n), g(n))) \\
mentioned by 2 result(s)

\smallskip

\noindent\textbf{\texttt{dom}}\quad\textit{operator}\\
a kernel term-former \\
mentioned by 18 result(s)

\smallskip

\noindent\textbf{\texttt{dual-norm(m, f)}}\quad\textit{functoid}\\
= iota(c\_, c\_ in rr and 0 <= c\_ and forall([x\_ in vec(m)], abs(f(x\_)) <= c\_ * (vnrm(m))(x\_)) and forall([d\_], d\_ ... \\
mentioned by 6 result(s)

\smallskip

\noindent\textbf{\texttt{dual-norm-on(m, s, f)}}\quad\textit{functoid}\\
= iota(c\_, c\_ in rr and 0 <= c\_ and forall([x\_ in s], abs(f(x\_)) <= c\_ * (vnrm(m))(x\_)) and forall([d\_], d\_ in r ... \\
mentioned by 8 result(s)

\smallskip


\section*{E}
\addcontentsline{toc}{section}{E}

\noindent\textbf{\texttt{elem-f(a, n, k, l)}}\quad\textit{functoid}\\
= matof(n, n, vnb-lambda([i, j], cartesian(interval(1, n), interval(1, n)), if(i = k and j = l or i = l and j =  ... \\
mentioned by 42 result(s)

\smallskip

\noindent\textbf{\texttt{elem-g(a, n, r, k, l)}}\quad\textit{functoid}\\
= matof(n, n, vnb-lambda([i, j], cartesian(interval(1, n), interval(1, n)), if(i = j, one(a), if(i = k and j = l ... \\
mentioned by 40 result(s)

\smallskip

\noindent\textbf{\texttt{elem-h(a, n, r, k)}}\quad\textit{functoid}\\
= matof(n, n, vnb-lambda([i, j], cartesian(interval(1, n), interval(1, n)), if(i = j, if(i = k, r, one(a)), zero ... \\
mentioned by 17 result(s)

\smallskip

\noindent\textbf{\texttt{embed(m)}}\quad\textit{functoid}\\
= vnb-lambda(u, pts(m), class(cauchy-setoid(m), embed-seq(m, u))) \\
mentioned by 3 result(s)

\smallskip

\noindent\textbf{\texttt{embed-seq(m, u)}}\quad\textit{functoid}\\
= vnb-lambda(n, nn, u)

\smallskip

\noindent\textbf{\texttt{entry(m, i, j)}}\quad\textit{functoid}\\
= nth(j, nth(i, m)) \\
mentioned by 230 result(s)

\smallskip

\noindent\textbf{\texttt{enum-fam(ag, f, phi, n)}}\quad\textit{functoid}\\
= vnb-lambda(i, nn, if(i in ord-segment(n), f(phi(i)), iden(ag))) \\
mentioned by 5 result(s)

\smallskip

\noindent\textbf{\texttt{eplus}}\quad\textit{operator}\\
a kernel term-former \\
mentioned by 6 result(s)

\smallskip

\noindent\textbf{\texttt{esum}}\quad\textit{operator}\\
a kernel term-former \\
mentioned by 7 result(s)

\smallskip

\noindent\textbf{\texttt{esup}}\quad\textit{operator}\\
a kernel term-former \\
mentioned by 5 result(s)

\smallskip

\noindent\textbf{\texttt{euclidean-gauges(s)}}\quad\textit{functoid}\\
= \{dg in fun(carr(s), nn): has-div-remainder(s, dg)\} \\
mentioned by 3 result(s)

\smallskip

\noindent\textbf{\texttt{euclidean-ring}}\quad\textit{refinement}\\
predicate \texttt{is-euclidean-ring}; same shape as integral-domain

\smallskip

\noindent\textbf{\texttt{ew}}\quad\textit{tactic}\\
\texttt{(ew term)} -- Existential witness: discharge a FORSOME goal by supplying the witness term. \\
declared in \texttt{tactics-help.scm}

\smallskip

\noindent\textbf{\texttt{exp}}\quad\textit{operator}\\
a kernel term-former

\smallskip

\noindent\textbf{\texttt{extends-on(s, g, f)}}\quad\textit{predicate}\\
reads: \$2 agrees with \$3 on \$1 \\
defined by \texttt{extends-on} \\
mentioned by 9 result(s)

\smallskip


\section*{F}
\addcontentsline{toc}{section}{F}

\noindent\textbf{\texttt{fact}}\quad\textit{tactic}\\
\texttt{(fact 'thm term ...)} -- Forward APPLICATION of a theorem: bring it in, instantiate its leading universals with the terms, and auto-detach every antecedent already in context, landing the consequent as a hypothesis. \\
declared in \texttt{tactics-help.scm}

\smallskip

\noindent\textbf{\texttt{factorial}}\quad\textit{defined-fn}\\
defined by \texttt{factorial-zero} \texttt{factorial-succ} \\
mentioned by 20 result(s)

\smallskip

\noindent\textbf{\texttt{falling}}\quad\textit{defined-fn}\\
defined by \texttt{falling-zero} \texttt{falling-succ} \\
mentioned by 9 result(s)

\smallskip

\noindent\textbf{\texttt{fam-of-list}}\quad\textit{defined-fn}\\
defined by \texttt{fam-of-list-apply} \\
mentioned by 18 result(s)

\smallskip

\noindent\textbf{\texttt{field}}\quad\textit{structure}\\
predicate \texttt{is-field}; slots: carr add mul neg zero one non-zero recip

\smallskip

\noindent\textbf{\texttt{field-additive-ag(s)}}\quad\textit{view}\\
a field viewed as an abelian-group \\
mentioned by 3 result(s)

\smallskip

\noindent\textbf{\texttt{field-as-euclidean-ring(s)}}\quad\textit{view}\\
a field viewed as an euclidean-ring \\
mentioned by 4 result(s)

\smallskip

\noindent\textbf{\texttt{field-as-integral-domain(s)}}\quad\textit{view}\\
a field viewed as an integral-domain \\
mentioned by 2 result(s)

\smallskip

\noindent\textbf{\texttt{field-multiplicative-group(s)}}\quad\textit{view}\\
a field viewed as a group \\
mentioned by 1 result(s)

\smallskip

\noindent\textbf{\texttt{fin-enum(s)}}\quad\textit{functoid}\\
= choice(bijection(ord-segment(card(s)), s)) \\
mentioned by 1 result(s)

\smallskip

\noindent\textbf{\texttt{finprod(m, f, s)}}\quad\textit{functoid}\\
= finsum(m, f, s)

\smallskip

\noindent\textbf{\texttt{finsum(ag, f, s)}}\quad\textit{functoid}\\
= sum-ag(ag, enum-fam(ag, f, fin-enum(s), card(s)), card(s)) \\
mentioned by 145 result(s)

\smallskip

\noindent\textbf{\texttt{finsupp(a, m)}}\quad\textit{functoid}\\
reads: the finitely-supported functions \$2 -> \$1 \\
= \{f in fun(carr(m), carr(a)): card(supp(a, m, f)) in nn\} \\
mentioned by 8 result(s)

\smallskip

\noindent\textbf{\texttt{fnrm(s)}}\quad\textit{accessor}\\
a structure slot \\
mentioned by 13 result(s)

\smallskip

\noindent\textbf{\texttt{focus}}\quad\textit{tactic}\\
\texttt{(focus n)} -- Switch the focus to the n-th open goal (1-based). \\
declared in \texttt{tactics-help.scm}

\smallskip

\noindent\textbf{\texttt{focus-leaf!}}\quad\textit{tactic}\\
\texttt{(focus-leaf! substr)} -- Focus the frontier leaf whose goal contains substr (never trust auto-advance).  [proof-local] \\
declared in \texttt{tactics-help.scm}

\smallskip

\noindent\textbf{\texttt{fr-ball(m, fam, n, eps)}}\quad\textit{functoid}\\
= \{x in vec(m): forall([k], k in nn and k <= n implies (fam(k))(x) < eps)\} \\
mentioned by 2 result(s)

\smallskip

\noindent\textbf{\texttt{fr-cauchy(m, fam, seq)}}\quad\textit{predicate}\\
reads: \$3 is Cauchy in (\$1, \$2) \\
defined by \texttt{fr-cauchy} \\
mentioned by 4 result(s)

\smallskip

\noindent\textbf{\texttt{fr-conv(m, fam, seq, lim)}}\quad\textit{predicate}\\
reads: \$3 converges to \$4 in (\$1, \$2) \\
defined by \texttt{fr-conv} \\
mentioned by 8 result(s)

\smallskip

\noindent\textbf{\texttt{fun}}\quad\textit{operator}\\
a kernel term-former \\
mentioned by 506 result(s)

\smallskip


\section*{G}
\addcontentsline{toc}{section}{G}

\noindent\textbf{\texttt{gauge(s)}}\quad\textit{functoid}\\
= choice(euclidean-gauges(s)) \\
mentioned by 15 result(s)

\smallskip

\noindent\textbf{\texttt{generates(md, n, u)}}\quad\textit{predicate}\\
reads: the \$2 vectors \$3 generate \$1 \\
defined by \texttt{generates} \\
mentioned by 8 result(s)

\smallskip

\noindent\textbf{\texttt{goal-status}}\quad\textit{tactic}\\
\texttt{(goal-status)} -- One-line summary: done / N open goals. \\
declared in \texttt{tactics-help.scm}

\smallskip

\noindent\textbf{\texttt{good-sub(m, s, f, t)}}\quad\textit{predicate}\\
reads: \$3 has a norm-preserving extension from \$2 to \$4, in \$1 \\
defined by \texttt{good-sub} \\
mentioned by 7 result(s)

\smallskip

\noindent\textbf{\texttt{greedy-chain(phi, grd, porel)}}\quad\textit{functoid}\\
= zkept(phi, grd, porel, card(grd)) \\
mentioned by 1 result(s)

\smallskip

\noindent\textbf{\texttt{grind}}\quad\textit{tactic}\\
\texttt{(grind)} -- Saturate the no-choice moves at the focus: decompose the goal connective (di) and break open the hypotheses (mac-h*), repeatedly, until neither fires.  The whole introduce-everything prefix in one step. \\
declared in \texttt{tactics-help.scm}

\smallskip

\noindent\textbf{\texttt{group}}\quad\textit{structure}\\
predicate \texttt{is-group}; slots: carr opr iden inv

\smallskip


\section*{H}
\addcontentsline{toc}{section}{H}

\noindent\textbf{\texttt{has-closed-graph(m1, fam1, m2, fam2, tt)}}\quad\textit{predicate}\\
reads: \$5 has a closed graph from (\$1, \$2) to (\$3, \$4) \\
defined by \texttt{has-closed-graph} \\
mentioned by 3 result(s)

\smallskip

\noindent\textbf{\texttt{has-div-remainder(s, deg)}}\quad\textit{predicate}\\
reads: \$1 has division with remainder for the degree function \$2 \\
defined by \texttt{has-div-remainder} \\
mentioned by 8 result(s)

\smallskip

\noindent\textbf{\texttt{has-fip(s, c)}}\quad\textit{predicate}\\
reads: \$2 has the finite intersection property in \$1 \\
defined by \texttt{has-fip} \\
mentioned by 4 result(s)

\smallskip

\noindent\textbf{\texttt{has-inverses(op, unit, invop, crr)}}\quad\textit{predicate}\\
reads: every element of \$4 has an inverse under \$1, given by \$3, with unit \$2 \\
defined by \texttt{has-inverses} \\
mentioned by 20 result(s)

\smallskip

\noindent\textbf{\texttt{have!}}\quad\textit{tactic}\\
\texttt{(have! CLAIM [THUNK])} -- Assert CLAIM as an intermediate step and carry on with it as a hypothesis: cut CLAIM, discharge the side goal from context (or by THUNK), and stay on the main branch. \\
declared in \texttt{tactics-help.scm}

\smallskip


\section*{I}
\addcontentsline{toc}{section}{I}

\noindent\textbf{\texttt{iden(s)}}\quad\textit{accessor}\\
a structure slot \\
mentioned by 102 result(s)

\smallskip

\noindent\textbf{\texttt{identmat(a, n)}}\quad\textit{functoid}\\
= matof(n, n, vnb-lambda([i, j], cartesian(interval(1, n), interval(1, n)), if(i = j, one(a), zero(a)))) \\
mentioned by 27 result(s)

\smallskip

\noindent\textbf{\texttt{idl(s)}}\quad\textit{accessor}\\
a structure slot \\
mentioned by 12 result(s)

\smallskip

\noindent\textbf{\texttt{ie}}\quad\textit{tactic}\\
\texttt{(ie hyp k)} -- Intersection elim: extract the k-th branch of an INTERSECTION-membership assumption. \\
declared in \texttt{tactics-help.scm}

\smallskip

\noindent\textbf{\texttt{if}}\quad\textit{operator}\\
a kernel term-former \\
mentioned by 46 result(s)

\smallskip

\noindent\textbf{\texttt{if-false}}\quad\textit{tactic}\\
\texttt{(if-false t)} -- Reduce an (IF p a b) term on the not-p branch: spawns (NOT p); the continuation gains (= (IF p a b) b). \\
declared in \texttt{tactics-help.scm}

\smallskip

\noindent\textbf{\texttt{if-true}}\quad\textit{tactic}\\
\texttt{(if-true t)} -- Reduce an (IF p a b) term on the p branch: spawns p as a subgoal; the continuation gains (= (IF p a b) a). \\
declared in \texttt{tactics-help.scm}

\smallskip

\noindent\textbf{\texttt{ii}}\quad\textit{tactic}\\
\texttt{(ii)} -- Intersection intro: prove (IN x (INTERSECTION ...)) for each branch. \\
declared in \texttt{tactics-help.scm}

\smallskip

\noindent\textbf{\texttt{imag-part}}\quad\textit{operator}\\
a kernel term-former

\smallskip

\noindent\textbf{\texttt{image}}\quad\textit{operator}\\
a kernel term-former \\
mentioned by 13 result(s)

\smallskip

\noindent\textbf{\texttt{in}}\quad\textit{primitive}\\
reads: \$1 is in \$2 \\
a kernel relation \\
mentioned by 2576 result(s)

\smallskip

\noindent\textbf{\texttt{ineq}}\quad\textit{tactic}\\
\texttt{(ineq i1 i2 ...)} -- Close a linear-inequality goal over RR (<= < > >= = between RR terms) as a consequence of the named assumptions (1-based indices), via the Fourier-Motzkin/Farkas oracle.  Linearizes over + - * and the binplus/binneg/bintimes aliases; every MAXIMAL non-arithmetic subterm is an atom that must be certified in RR.  (Does NOT see through a generic ring's (ADD s)/(MUL s) -- those become opaque atoms.) \\
declared in \texttt{tactics-help.scm}

\smallskip

\noindent\textbf{\texttt{inf-subsets(a)}}\quad\textit{functoid}\\
= \{s in power(a): not(card(s) in nn)\} \\
mentioned by 13 result(s)

\smallskip

\noindent\textbf{\texttt{injection}}\quad\textit{operator}\\
a kernel term-former \\
mentioned by 16 result(s)

\smallskip

\noindent\textbf{\texttt{injective*(f)}}\quad\textit{predicate}\\
reads: \$1 is injective \\
defined by \texttt{injective*} \\
mentioned by 3 result(s)

\smallskip

\noindent\textbf{\texttt{insert-last(phi, x, n)}}\quad\textit{functoid}\\
= vnb-lambda(i, nn, if(i in ord-segment(n), phi(i), x))

\smallskip

\noindent\textbf{\texttt{inst}}\quad\textit{tactic}\\
\texttt{(inst forall-hyp term)} -- Instantiate a universally-quantified assumption at `term', adding the instance to context. \\
declared in \texttt{tactics-help.scm}

\smallskip

\noindent\textbf{\texttt{inst+}}\quad\textit{tactic}\\
\texttt{(inst+ forall-hyp term)} -- Instantiate an in-context universal at `term', then forward-detach any guards whose antecedents are in context. \\
declared in \texttt{tactics-help.scm}

\smallskip

\noindent\textbf{\texttt{integral-domain}}\quad\textit{refinement}\\
predicate \texttt{is-integral-domain}; same shape as commutative-ring

\smallskip

\noindent\textbf{\texttt{interior(s, a)}}\quad\textit{functoid}\\
= \{x in pts(s): forsome([u], is-open(s, u) and x in u and u subset a)\} \\
mentioned by 2 result(s)

\smallskip

\noindent\textbf{\texttt{intersection}}\quad\textit{operator}\\
a kernel term-former \\
mentioned by 22 result(s)

\smallskip

\noindent\textbf{\texttt{interval(a, b)}}\quad\textit{functoid}\\
= \{i in nn: a <= i and i <= b\} \\
mentioned by 218 result(s)

\smallskip

\noindent\textbf{\texttt{inv(s)}}\quad\textit{accessor}\\
a structure slot \\
mentioned by 29 result(s)

\smallskip

\noindent\textbf{\texttt{inverse-bij(phi, x, y)}}\quad\textit{functoid}\\
= vnb-lambda(y\_, y, choice(\{x\_ in x: phi(x\_) = y\_\})) \\
mentioned by 6 result(s)

\smallskip

\noindent\textbf{\texttt{iota}}\quad\textit{operator}\\
a kernel term-former \\
mentioned by 2 result(s)

\smallskip

\noindent\textbf{\texttt{iota-d}}\quad\textit{tactic}\\
\texttt{(iota-d term)} -- Definite-description: discharge the IOTA uniqueness obligation for `term'. \\
declared in \texttt{tactics-help.scm}

\smallskip

\noindent\textbf{\texttt{is-abelian-group(s)}}\quad\textit{predicate}\\
reads: an abelian group \\
a predicate \\
mentioned by 106 result(s)

\smallskip

\noindent\textbf{\texttt{is-absolutely-summable(grp, f)}}\quad\textit{predicate}\\
reads: \$2 is absolutely summable in \$1 \\
defined by \texttt{is-absolutely-summable} \\
mentioned by 3 result(s)

\smallskip

\noindent\textbf{\texttt{is-algebra-of-sets(omega, ca)}}\quad\textit{predicate}\\
reads: \$2 is an algebra of subsets of \$1 \\
defined by \texttt{is-algebra-of-sets} \\
mentioned by 2 result(s)

\smallskip

\noindent\textbf{\texttt{is-associative(op, crr)}}\quad\textit{predicate}\\
reads: \$1 is associative on \$2 \\
defined by \texttt{is-associative} \\
mentioned by 26 result(s)

\smallskip

\noindent\textbf{\texttt{is-bounded-linear-functional(m, f)}}\quad\textit{predicate}\\
reads: \$2 is a bounded linear functional on \$1 \\
defined by \texttt{is-bounded-linear-functional} \\
mentioned by 15 result(s)

\smallskip

\noindent\textbf{\texttt{is-bounded-linear-functional-on(m, s, f)}}\quad\textit{predicate}\\
reads: \$3 is a bounded linear functional on the subspace \$2 of \$1 \\
defined by \texttt{is-bounded-linear-functional-on} \\
mentioned by 10 result(s)

\smallskip

\noindent\textbf{\texttt{is-bounded-metric-space(s)}}\quad\textit{predicate}\\
reads: a bounded metric space \\
defined by \texttt{is-bounded-metric-space} \\
mentioned by 7 result(s)

\smallskip

\noindent\textbf{\texttt{is-cauchy-seq(s, f)}}\quad\textit{predicate}\\
reads: Cauchy \\
defined by \texttt{is-cauchy-seq} \\
mentioned by 11 result(s)

\smallskip

\noindent\textbf{\texttt{is-chain(grd, porel, ch)}}\quad\textit{predicate}\\
reads: \$3 is a chain in \$1 under \$2 \\
defined by \texttt{is-chain} \\
mentioned by 11 result(s)

\smallskip

\noindent\textbf{\texttt{is-closed(s, a)}}\quad\textit{predicate}\\
reads: \$2 is closed in \$1 \\
defined by \texttt{is-closed} \\
mentioned by 6 result(s)

\smallskip

\noindent\textbf{\texttt{is-comm-monoid(s)}}\quad\textit{predicate}\\
reads: a comm monoid \\
a predicate \\
mentioned by 32 result(s)

\smallskip

\noindent\textbf{\texttt{is-commutative(op, crr)}}\quad\textit{predicate}\\
reads: \$1 is commutative on \$2 \\
defined by \texttt{is-commutative} \\
mentioned by 20 result(s)

\smallskip

\noindent\textbf{\texttt{is-commutative-ring(s)}}\quad\textit{predicate}\\
reads: a commutative ring \\
a predicate \\
mentioned by 108 result(s)

\smallskip

\noindent\textbf{\texttt{is-compact(s)}}\quad\textit{predicate}\\
reads: \$1 is compact \\
defined by \texttt{is-compact} \\
mentioned by 18 result(s)

\smallskip

\noindent\textbf{\texttt{is-complete(s)}}\quad\textit{predicate}\\
reads: complete \\
defined by \texttt{is-complete} \\
mentioned by 12 result(s)

\smallskip

\noindent\textbf{\texttt{is-cont-lin(m1, fam1, m2, fam2, tt)}}\quad\textit{predicate}\\
reads: \$5 is a continuous linear map from (\$1, \$2) to (\$3, \$4) \\
defined by \texttt{is-cont-lin} \\
mentioned by 5 result(s)

\smallskip

\noindent\textbf{\texttt{is-continuous(s, t, f)}}\quad\textit{predicate}\\
reads: \$3 is continuous from \$1 to \$2 \\
defined by \texttt{is-continuous} \\
mentioned by 16 result(s)

\smallskip

\noindent\textbf{\texttt{is-continuous-at(s, t, f, a)}}\quad\textit{predicate}\\
reads: \$3 is continuous at \$4 \\
defined by \texttt{is-continuous-at} \\
mentioned by 45 result(s)

\smallskip

\noindent\textbf{\texttt{is-countable-pseudometric-family(fam, ground)}}\quad\textit{predicate}\\
reads: \$1 is a countable family of pseudometrics on \$2 \\
defined by \texttt{is-countable-pseudometric-family} \\
mentioned by 4 result(s)

\smallskip

\noindent\textbf{\texttt{is-dense-seq(s, dseq)}}\quad\textit{predicate}\\
reads: \$2 is a dense sequence in \$1 \\
defined by \texttt{is-dense-seq} \\
mentioned by 4 result(s)

\smallskip

\noindent\textbf{\texttt{is-diagonal(a, m, n, d)}}\quad\textit{predicate}\\
reads: \$4 is a diagonal \$2-by-\$3 matrix over \$1 \\
defined by \texttt{is-diagonal} \\
mentioned by 7 result(s)

\smallskip

\noindent\textbf{\texttt{is-diff-at(f, a, l)}}\quad\textit{predicate}\\
reads: \$1 is differentiable at \$2, with derivative \$3 \\
defined by \texttt{is-diff-at} \\
mentioned by 36 result(s)

\smallskip

\noindent\textbf{\texttt{is-diff-at-v(m, f, a, l)}}\quad\textit{predicate}\\
reads: \$2 is differentiable at \$3 with derivative \$4, as a curve in \$1 \\
defined by \texttt{is-diff-at-v} \\
mentioned by 4 result(s)

\smallskip

\noindent\textbf{\texttt{is-distributive(addop, mulop, crr)}}\quad\textit{predicate}\\
reads: \$2 distributes over \$1 on \$3 \\
defined by \texttt{is-distributive} \\
mentioned by 10 result(s)

\smallskip

\noindent\textbf{\texttt{is-eps-cauchy-seq(s, eps, y)}}\quad\textit{predicate}\\
reads: \$3 is \$2-Cauchy in \$1 \\
defined by \texttt{is-eps-cauchy-seq} \\
mentioned by 3 result(s)

\smallskip

\noindent\textbf{\texttt{is-equicontinuous(s, t, fam)}}\quad\textit{predicate}\\
reads: \$3 is equicontinuous from \$1 to \$2 \\
defined by \texttt{is-equicontinuous} \\
mentioned by 5 result(s)

\smallskip

\noindent\textbf{\texttt{is-equivalence(rho, crr)}}\quad\textit{predicate}\\
reads: \$1 is an equivalence relation on \$2 \\
defined by \texttt{is-equivalence} \\
mentioned by 5 result(s)

\smallskip

\noindent\textbf{\texttt{is-euclidean-ring(s)}}\quad\textit{predicate}\\
reads: a Euclidean ring \\
a predicate \\
mentioned by 33 result(s)

\smallskip

\noindent\textbf{\texttt{is-field(s)}}\quad\textit{predicate}\\
reads: a field \\
a predicate \\
mentioned by 80 result(s)

\smallskip

\noindent\textbf{\texttt{is-finite-cover(c, a)}}\quad\textit{predicate}\\
reads: \$1 is a finite cover of \$2 \\
defined by \texttt{is-finite-cover} \\
mentioned by 5 result(s)

\smallskip

\noindent\textbf{\texttt{is-finite-dimensional(m)}}\quad\textit{predicate}\\
reads: \$1 is finite dimensional \\
defined by \texttt{is-finite-dimensional} \\
mentioned by 8 result(s)

\smallskip

\noindent\textbf{\texttt{is-frechet-structure(m, fam)}}\quad\textit{predicate}\\
reads: (\$1, \$2) is a Frechet space \\
defined by \texttt{is-frechet-structure} \\
mentioned by 5 result(s)

\smallskip

\noindent\textbf{\texttt{is-fun}}\quad\textit{predicate}\\
a predicate \\
mentioned by 2 result(s)

\smallskip

\noindent\textbf{\texttt{is-gauge-countable(s)}}\quad\textit{predicate}\\
reads: \$1 has a countable pseudometric gauge \\
defined by \texttt{is-gauge-countable} \\
mentioned by 7 result(s)

\smallskip

\noindent\textbf{\texttt{is-group(s)}}\quad\textit{predicate}\\
reads: a group \\
a predicate \\
mentioned by 20 result(s)

\smallskip

\noindent\textbf{\texttt{is-group-norm(nm, op, invop, unit, crr)}}\quad\textit{predicate}\\
reads: \$1 is a group norm for \$2 on \$5 \\
defined by \texttt{is-group-norm} \\
mentioned by 4 result(s)

\smallskip

\noindent\textbf{\texttt{is-hausdorff(s)}}\quad\textit{predicate}\\
reads: \$1 is Hausdorff \\
defined by \texttt{is-hausdorff} \\
mentioned by 6 result(s)

\smallskip

\noindent\textbf{\texttt{is-hom-abelian-group(a, b, f)}}\quad\textit{predicate}\\
reads: \$3 is a homomorphism from \$1 to \$2 \\
a predicate \\
mentioned by 9 result(s)

\smallskip

\noindent\textbf{\texttt{is-hom-comm-monoid(a, b, f)}}\quad\textit{predicate}\\
reads: \$3 is a homomorphism from \$1 to \$2 \\
a predicate \\
mentioned by 3 result(s)

\smallskip

\noindent\textbf{\texttt{is-hom-commutative-ring(a, b, f)}}\quad\textit{predicate}\\
reads: \$3 is a homomorphism from \$1 to \$2 \\
a predicate \\
mentioned by 7 result(s)

\smallskip

\noindent\textbf{\texttt{is-hom-euclidean-ring(a, b, f)}}\quad\textit{predicate}\\
reads: \$3 is a homomorphism from \$1 to \$2 \\
a predicate \\
mentioned by 3 result(s)

\smallskip

\noindent\textbf{\texttt{is-hom-field(a, b, f)}}\quad\textit{predicate}\\
reads: \$3 is a homomorphism from \$1 to \$2 \\
a predicate \\
mentioned by 5 result(s)

\smallskip

\noindent\textbf{\texttt{is-hom-group(a, b, f)}}\quad\textit{predicate}\\
reads: \$3 is a homomorphism from \$1 to \$2 \\
a predicate \\
mentioned by 2 result(s)

\smallskip

\noindent\textbf{\texttt{is-hom-integral-domain(a, b, f)}}\quad\textit{predicate}\\
reads: \$3 is a homomorphism from \$1 to \$2 \\
a predicate \\
mentioned by 8 result(s)

\smallskip

\noindent\textbf{\texttt{is-hom-metric-space(a, b, f)}}\quad\textit{predicate}\\
reads: \$3 is an isometry from \$1 to \$2 \\
a predicate \\
mentioned by 4 result(s)

\smallskip

\noindent\textbf{\texttt{is-hom-metrizable-top-space(s, t, f)}}\quad\textit{predicate}\\
reads: \$3 is continuous from \$1 to \$2 \\
a predicate \\
mentioned by 3 result(s)

\smallskip

\noindent\textbf{\texttt{is-hom-module(a, b, f)}}\quad\textit{predicate}\\
reads: \$3 is a homomorphism from \$1 to \$2 \\
a predicate \\
mentioned by 6 result(s)

\smallskip

\noindent\textbf{\texttt{is-hom-monoid(a, b, f)}}\quad\textit{predicate}\\
reads: \$3 is a homomorphism from \$1 to \$2 \\
a predicate \\
mentioned by 4 result(s)

\smallskip

\noindent\textbf{\texttt{is-hom-normed-ag(a, b, f)}}\quad\textit{predicate}\\
reads: \$3 is a homomorphism from \$1 to \$2 \\
a predicate \\
mentioned by 4 result(s)

\smallskip

\noindent\textbf{\texttt{is-hom-normed-field(a, b, f)}}\quad\textit{predicate}\\
reads: \$3 is a homomorphism from \$1 to \$2 \\
a predicate \\
mentioned by 5 result(s)

\smallskip

\noindent\textbf{\texttt{is-hom-normed-vector-space(a, b, f)}}\quad\textit{predicate}\\
reads: \$3 is a homomorphism from \$1 to \$2 \\
a predicate \\
mentioned by 4 result(s)

\smallskip

\noindent\textbf{\texttt{is-hom-pid(a, b, f)}}\quad\textit{predicate}\\
reads: \$3 is a homomorphism from \$1 to \$2 \\
a predicate \\
mentioned by 2 result(s)

\smallskip

\noindent\textbf{\texttt{is-hom-pseudometric-space(a, b, f)}}\quad\textit{predicate}\\
reads: \$3 is a homomorphism from \$1 to \$2 \\
a predicate \\
mentioned by 2 result(s)

\smallskip

\noindent\textbf{\texttt{is-hom-ring(a, b, f)}}\quad\textit{predicate}\\
reads: \$3 is a homomorphism from \$1 to \$2 \\
a predicate \\
mentioned by 8 result(s)

\smallskip

\noindent\textbf{\texttt{is-hom-ringoid(a, b, f)}}\quad\textit{predicate}\\
reads: \$3 is a ringoid homomorphism from \$1 to \$2 \\
a predicate \\
mentioned by 2 result(s)

\smallskip

\noindent\textbf{\texttt{is-hom-semigroup(a, b, f)}}\quad\textit{predicate}\\
reads: \$3 is a homomorphism from \$1 to \$2 \\
a predicate \\
mentioned by 2 result(s)

\smallskip

\noindent\textbf{\texttt{is-hom-setoid(a, b, f1, f2)}}\quad\textit{predicate}\\
reads: \$3 is a setoid homomorphism from \$1 to \$2 \\
a predicate \\
mentioned by 2 result(s)

\smallskip

\noindent\textbf{\texttt{is-hom-top-space(s, t, f)}}\quad\textit{predicate}\\
reads: \$3 is continuous from \$1 to \$2 \\
a predicate \\
mentioned by 2 result(s)

\smallskip

\noindent\textbf{\texttt{is-hom-vector-space(a, b, f)}}\quad\textit{predicate}\\
reads: \$3 is a homomorphism from \$1 to \$2 \\
a predicate \\
mentioned by 2 result(s)

\smallskip

\noindent\textbf{\texttt{is-ideal(s, i)}}\quad\textit{predicate}\\
reads: \$2 is an ideal of \$1 \\
defined by \texttt{is-ideal} \\
mentioned by 9 result(s)

\smallskip

\noindent\textbf{\texttt{is-ideal-in(i\_, carr\_, add\_, mul\_, neg\_, zero\_)}}\quad\textit{predicate}\\
reads: \$1 is a two-sided ideal of the ring with carrier \$2 \\
defined by \texttt{is-ideal-in} \\
mentioned by 4 result(s)

\smallskip

\noindent\textbf{\texttt{is-identity(op, unit, crr)}}\quad\textit{predicate}\\
reads: \$2 is an identity for \$1 on \$3 \\
defined by \texttt{is-identity} \\
mentioned by 24 result(s)

\smallskip

\noindent\textbf{\texttt{is-integral-domain(s)}}\quad\textit{predicate}\\
reads: an integral domain \\
a predicate \\
mentioned by 19 result(s)

\smallskip

\noindent\textbf{\texttt{is-invertible-mat(a, n, u)}}\quad\textit{predicate}\\
reads: \$3 is an invertible \$2-by-\$2 matrix over \$1 \\
defined by \texttt{is-invertible-mat} \\
mentioned by 17 result(s)

\smallskip

\noindent\textbf{\texttt{is-isometry(s, t, f)}}\quad\textit{predicate}\\
reads: \$3 is an isometry from \$1 to \$2 \\
defined by \texttt{is-isometry} \\
mentioned by 2 result(s)

\smallskip

\noindent\textbf{\texttt{is-k-linear(m, f)}}\quad\textit{predicate}\\
reads: \$2 is a K-linear functional on \$1 \\
defined by \texttt{is-k-linear} \\
mentioned by 3 result(s)

\smallskip

\noindent\textbf{\texttt{is-k-linear-on(m, s, f)}}\quad\textit{predicate}\\
reads: \$3 is a K-linear functional on the submodule \$2 of \$1 \\
defined by \texttt{is-k-linear-on} \\
mentioned by 3 result(s)

\smallskip

\noindent\textbf{\texttt{is-linear-functional(m, f)}}\quad\textit{predicate}\\
reads: \$2 is a linear functional on \$1 \\
defined by \texttt{is-linear-functional} \\
mentioned by 6 result(s)

\smallskip

\noindent\textbf{\texttt{is-linear-functional-on(m, s, f)}}\quad\textit{predicate}\\
reads: \$3 is a linear functional on the subspace \$2 of \$1 \\
defined by \texttt{is-linear-functional-on} \\
mentioned by 12 result(s)

\smallskip

\noindent\textbf{\texttt{is-linear-map(m1, m2, tt)}}\quad\textit{predicate}\\
reads: \$3 is a linear map from \$1 to \$2 \\
defined by \texttt{is-linear-map} \\
mentioned by 9 result(s)

\smallskip

\noindent\textbf{\texttt{is-maximal(grd, porel, mx)}}\quad\textit{predicate}\\
reads: \$3 is maximal in \$1 under \$2 \\
defined by \texttt{is-maximal} \\
mentioned by 4 result(s)

\smallskip

\noindent\textbf{\texttt{is-meager(s, a)}}\quad\textit{predicate}\\
reads: \$2 is meager in \$1 \\
defined by \texttt{is-meager} \\
mentioned by 4 result(s)

\smallskip

\noindent\textbf{\texttt{is-metric(dst, crr)}}\quad\textit{predicate}\\
reads: \$1 is a metric on \$2 \\
defined by \texttt{is-metric} \\
mentioned by 4 result(s)

\smallskip

\noindent\textbf{\texttt{is-metric-space(s)}}\quad\textit{predicate}\\
reads: a metric space \\
a predicate \\
mentioned by 125 result(s)

\smallskip

\noindent\textbf{\texttt{is-metrizable-top-space(s)}}\quad\textit{predicate}\\
reads: \$1 is a metrizable topological space \\
a predicate \\
mentioned by 13 result(s)

\smallskip

\noindent\textbf{\texttt{is-module(s)}}\quad\textit{predicate}\\
reads: a module \\
a predicate \\
mentioned by 146 result(s)

\smallskip

\noindent\textbf{\texttt{is-monoid(s)}}\quad\textit{predicate}\\
reads: a monoid \\
a predicate \\
mentioned by 45 result(s)

\smallskip

\noindent\textbf{\texttt{is-ms-sequence(ms)}}\quad\textit{predicate}\\
reads: \$1 is a sequence of metric spaces \\
defined by \texttt{is-ms-sequence} \\
mentioned by 12 result(s)

\smallskip

\noindent\textbf{\texttt{is-noetherian(m)}}\quad\textit{predicate}\\
reads: \$1 is Noetherian \\
defined by \texttt{is-noetherian} \\
mentioned by 5 result(s)

\smallskip

\noindent\textbf{\texttt{is-nonmeager(s, a)}}\quad\textit{predicate}\\
reads: \$2 is nonmeager in \$1 \\
defined by \texttt{is-nonmeager} \\
mentioned by 3 result(s)

\smallskip

\noindent\textbf{\texttt{is-norm(nm, addop, mulop, zr, crr)}}\quad\textit{predicate}\\
reads: \$1 is a norm on \$5 \\
defined by \texttt{is-norm} \\
mentioned by 4 result(s)

\smallskip

\noindent\textbf{\texttt{is-normed-ag(s)}}\quad\textit{predicate}\\
reads: a normed ag \\
a predicate \\
mentioned by 61 result(s)

\smallskip

\noindent\textbf{\texttt{is-normed-field(s)}}\quad\textit{predicate}\\
reads: a normed field \\
a predicate \\
mentioned by 161 result(s)

\smallskip

\noindent\textbf{\texttt{is-normed-vector-space(s)}}\quad\textit{predicate}\\
reads: a normed vector space \\
a predicate \\
mentioned by 231 result(s)

\smallskip

\noindent\textbf{\texttt{is-nowhere-dense(s, a)}}\quad\textit{predicate}\\
reads: \$2 is nowhere dense in \$1 \\
defined by \texttt{is-nowhere-dense} \\
mentioned by 4 result(s)

\smallskip

\noindent\textbf{\texttt{is-open(s, u)}}\quad\textit{predicate}\\
reads: \$2 is open in \$1 \\
defined by \texttt{is-open} \\
mentioned by 15 result(s)

\smallskip

\noindent\textbf{\texttt{is-open-cover(s, c)}}\quad\textit{predicate}\\
reads: \$2 is an open cover of \$1 \\
defined by \texttt{is-open-cover} \\
mentioned by 7 result(s)

\smallskip

\noindent\textbf{\texttt{is-open-gauge(fam, ground, u)}}\quad\textit{predicate}\\
reads: \$3 is open in the gauge topology of the family \$1 on \$2 \\
defined by \texttt{is-open-gauge} \\
mentioned by 2 result(s)

\smallskip

\noindent\textbf{\texttt{is-open-lin-map(m1, fam1, m2, fam2, tt)}}\quad\textit{predicate}\\
reads: \$5 is an open linear map from (\$1, \$2) to (\$3, \$4) \\
defined by \texttt{is-open-lin-map} \\
mentioned by 3 result(s)

\smallskip

\noindent\textbf{\texttt{is-ord(x)}}\quad\textit{predicate}\\
reads: \$1 is an ordinal \\
a predicate \\
mentioned by 2 result(s)

\smallskip

\noindent\textbf{\texttt{is-partial-order(grd, porel)}}\quad\textit{predicate}\\
reads: \$2 partially orders \$1 \\
defined by \texttt{is-partial-order} \\
mentioned by 13 result(s)

\smallskip

\noindent\textbf{\texttt{is-pid(s)}}\quad\textit{predicate}\\
reads: a pid \\
a predicate \\
mentioned by 6 result(s)

\smallskip

\noindent\textbf{\texttt{is-pseudometric(dst, crr)}}\quad\textit{predicate}\\
reads: \$1 is a pseudometric on \$2 \\
defined by \texttt{is-pseudometric} \\
mentioned by 4 result(s)

\smallskip

\noindent\textbf{\texttt{is-pseudometric-space(s)}}\quad\textit{predicate}\\
reads: a pseudometric space \\
a predicate \\
mentioned by 8 result(s)

\smallskip

\noindent\textbf{\texttt{is-r-net(s, f, a, r)}}\quad\textit{predicate}\\
reads: \$2 is an \$4-net for \$3 in \$1 \\
defined by \texttt{is-r-net} \\
mentioned by 5 result(s)

\smallskip

\noindent\textbf{\texttt{is-ring(s)}}\quad\textit{predicate}\\
reads: a ring \\
a predicate \\
mentioned by 387 result(s)

\smallskip

\noindent\textbf{\texttt{is-ringoid(s)}}\quad\textit{predicate}\\
reads: a ringoid \\
a predicate \\
mentioned by 78 result(s)

\smallskip

\noindent\textbf{\texttt{is-semigroup(s)}}\quad\textit{predicate}\\
reads: a semigroup \\
a predicate \\
mentioned by 8 result(s)

\smallskip

\noindent\textbf{\texttt{is-seminorm(m, p)}}\quad\textit{predicate}\\
reads: \$2 is a seminorm on \$1 \\
defined by \texttt{is-seminorm} \\
mentioned by 5 result(s)

\smallskip

\noindent\textbf{\texttt{is-seminorm-family(m, fam)}}\quad\textit{predicate}\\
reads: \$2 is a separating countable family of seminorms on \$1 \\
defined by \texttt{is-seminorm-family} \\
mentioned by 4 result(s)

\smallskip

\noindent\textbf{\texttt{is-separable(s)}}\quad\textit{predicate}\\
reads: \$1 is separable \\
defined by \texttt{is-separable} \\
mentioned by 3 result(s)

\smallskip

\noindent\textbf{\texttt{is-set(x)}}\quad\textit{predicate}\\
reads: \$1 is a set \\
a predicate \\
mentioned by 2 result(s)

\smallskip

\noindent\textbf{\texttt{is-setoid(s)}}\quad\textit{predicate}\\
reads: a setoid \\
a predicate \\
mentioned by 25 result(s)

\smallskip

\noindent\textbf{\texttt{is-sigma-algebra(omega, ca)}}\quad\textit{predicate}\\
reads: \$2 is a sigma-algebra of subsets of \$1 \\
defined by \texttt{is-sigma-algebra} \\
mentioned by 2 result(s)

\smallskip

\noindent\textbf{\texttt{is-strictly-below(porel, x, y)}}\quad\textit{predicate}\\
reads: \$2 is strictly below \$3 under \$1 \\
defined by \texttt{is-strictly-below} \\
mentioned by 11 result(s)

\smallskip

\noindent\textbf{\texttt{is-submodule(m, s)}}\quad\textit{predicate}\\
reads: \$2 is a submodule of \$1 \\
defined by \texttt{is-submodule} \\
mentioned by 39 result(s)

\smallskip

\noindent\textbf{\texttt{is-subsequence(s, y, f)}}\quad\textit{predicate}\\
reads: \$2 is a subsequence of \$3 in \$1 \\
defined by \texttt{is-subsequence} \\
mentioned by 2 result(s)

\smallskip

\noindent\textbf{\texttt{is-subspace(m, s)}}\quad\textit{predicate}\\
reads: \$2 is a subspace of \$1 \\
defined by \texttt{is-subspace} \\
mentioned by 2 result(s)

\smallskip

\noindent\textbf{\texttt{is-summable(grp, f)}}\quad\textit{predicate}\\
reads: \$2 is summable in \$1 \\
defined by \texttt{is-summable} \\
mentioned by 3 result(s)

\smallskip

\noindent\textbf{\texttt{is-top-space(s)}}\quad\textit{predicate}\\
reads: \$1 is a topological space \\
a predicate \\
mentioned by 19 result(s)

\smallskip

\noindent\textbf{\texttt{is-unif-cauchy(s, fam)}}\quad\textit{predicate}\\
reads: \$2 is uniformly Cauchy on \$1 \\
defined by \texttt{is-unif-cauchy} \\
mentioned by 4 result(s)

\smallskip

\noindent\textbf{\texttt{is-uniformly-continuous(s, t, f)}}\quad\textit{predicate}\\
reads: \$3 is uniformly continuous from \$1 to \$2 \\
defined by \texttt{is-uniformly-continuous} \\
mentioned by 3 result(s)

\smallskip

\noindent\textbf{\texttt{is-upper-bound(grd, porel, ch, b)}}\quad\textit{predicate}\\
reads: \$4 is an upper bound of \$3 in \$1 under \$2 \\
defined by \texttt{is-upper-bound} \\
mentioned by 6 result(s)

\smallskip

\noindent\textbf{\texttt{is-vector-space(s)}}\quad\textit{predicate}\\
reads: a vector space \\
a predicate \\
mentioned by 9 result(s)

\smallskip


\section*{K}
\addcontentsline{toc}{section}{K}

\noindent\textbf{\texttt{keep-set(phi, grd, porel, kset, alpha)}}\quad\textit{functoid}\\
= \{y in grd: y = phi(alpha) and forall([z in kset], [z, y] in porel)\} \\
mentioned by 5 result(s)

\smallskip


\section*{L}
\addcontentsline{toc}{section}{L}

\noindent\textbf{\texttt{lam-b}}\quad\textit{tactic}\\
\texttt{(lam-b)} -- VNB-LAMBDA beta: reduce an applied lambda to its substituted body. \\
declared in \texttt{tactics-help.scm}

\smallskip

\noindent\textbf{\texttt{lam-b-h}}\quad\textit{tactic}\\
\texttt{(lam-b-h hyp)} -- VNB-LAMBDA beta in a cited ASSUMPTION -- what mac-h is to mac. \\
declared in \texttt{tactics-help.scm}

\smallskip

\noindent\textbf{\texttt{lam-t}}\quad\textit{tactic}\\
\texttt{(lam-t)} -- VNB-LAMBDA typing: reduce (IN (VNB-LAMBDA ...) (FUN A B)) to TWO obligations -- the body, and that the domain is a set. \\
declared in \texttt{tactics-help.scm}

\smallskip

\noindent\textbf{\texttt{lastcoeff-set(md, p, u, sm)}}\quad\textit{functoid}\\
= \{r\_ in carr(scal(md)): forsome([c\_ in mat(1, succ(p), carr(scal(md)))], entry(c\_, 1, succ(p)) = r\_ and entry(m ... \\
mentioned by 3 result(s)

\smallskip

\noindent\textbf{\texttt{len-r}}\quad\textit{tactic}\\
\texttt{(len-r)} -- Reduce a LENGTH applied to a literal LIST in the goal. \\
declared in \texttt{tactics-help.scm}

\smallskip

\noindent\textbf{\texttt{length}}\quad\textit{operator}\\
a kernel term-former \\
mentioned by 42 result(s)

\smallskip

\noindent\textbf{\texttt{limit-ord}}\quad\textit{operator}\\
a kernel term-former \\
mentioned by 11 result(s)

\smallskip

\noindent\textbf{\texttt{line(m, v)}}\quad\textit{functoid}\\
= \{y\_ in vec(m): forsome([r\_ in rr], y\_ = (act(m))(r\_, v))\} \\
mentioned by 3 result(s)

\smallskip

\noindent\textbf{\texttt{list}}\quad\textit{operator}\\
a kernel term-former \\
mentioned by 111 result(s)

\smallskip

\noindent\textbf{\texttt{little-o-at(g, a)}}\quad\textit{predicate}\\
reads: \$1 is little-o at \$2 \\
defined by \texttt{little-o-at} \\
mentioned by 6 result(s)

\smallskip


\section*{M}
\addcontentsline{toc}{section}{M}

\noindent\textbf{\texttt{mac}}\quad\textit{tactic}\\
\texttt{(mac 'name)} -- Rewrite the GOAL with an equivalence macete (unfold a definition, apply an iff/=/== law).  Fires only where the macete's side-conditions already hold in context. \\
declared in \texttt{tactics-help.scm}

\smallskip

\noindent\textbf{\texttt{mac-h}}\quad\textit{tactic}\\
\texttt{(mac-h 'name hyp)} -- Rewrite a cited ASSUMPTION in place with an equivalence macete; any side-condition not already in context is spawned as a new subgoal. \\
declared in \texttt{tactics-help.scm}

\smallskip

\noindent\textbf{\texttt{mac-h*}}\quad\textit{tactic}\\
\texttt{(mac-h*)} -- Saturating mac-h: repeatedly unfold every defined predicate in the hypotheses and split the conjunctions they expose, until nothing is left folded.  One step in the trace instead of a dozen. \\
declared in \texttt{tactics-help.scm}

\smallskip

\noindent\textbf{\texttt{macm}}\quad\textit{tactic}\\
\texttt{(macm 'name)} -- Like mac, but a CONDITIONAL macete fires even when its side-conditions are not yet in context: each unmet condition is left as a new subgoal. \\
declared in \texttt{tactics-help.scm}

\smallskip

\noindent\textbf{\texttt{magnitude}}\quad\textit{operator}\\
a kernel term-former \\
mentioned by 15 result(s)

\smallskip

\noindent\textbf{\texttt{make-set}}\quad\textit{operator}\\
a kernel term-former \\
mentioned by 7 result(s)

\smallskip

\noindent\textbf{\texttt{mat(m, n, x)}}\quad\textit{functoid}\\
= \{p in matrix(x): size(p) = [m, n]\} \\
mentioned by 204 result(s)

\smallskip

\noindent\textbf{\texttt{mat-equiv(a, m, n, c, d)}}\quad\textit{predicate}\\
reads: \$4 and \$5 are equivalent \$2-by-\$3 matrices over \$1 \\
defined by \texttt{mat-equiv} \\
mentioned by 24 result(s)

\smallskip

\noindent\textbf{\texttt{mat-ring(a, n)}}\quad\textit{functoid}\\
= [mat(n, n, carr(a)), vnb-lambda([p, q], cartesian(mat(n, n, carr(a)), mat(n, n, carr(a))), matadd(a, p, q)), v ... \\
mentioned by 15 result(s)

\smallskip

\noindent\textbf{\texttt{matact(md, p, u)}}\quad\textit{functoid}\\
= matof(nth(1, size(p)), nth(2, size(u)), vnb-lambda([i, c], cartesian(interval(1, nth(1, size(p))), interval(1, ... \\
mentioned by 42 result(s)

\smallskip

\noindent\textbf{\texttt{matadd(a, p, q)}}\quad\textit{functoid}\\
= matof(nth(1, size(p)), nth(2, size(p)), vnb-lambda([i, j], cartesian(interval(1, nth(1, size(p))), interval(1, ... \\
mentioned by 24 result(s)

\smallskip

\noindent\textbf{\texttt{matmul(a, p, q)}}\quad\textit{functoid}\\
= matof(nth(1, size(p)), nth(2, size(q)), vnb-lambda([i, k], cartesian(interval(1, nth(1, size(p))), interval(1, ... \\
mentioned by 62 result(s)

\smallskip

\noindent\textbf{\texttt{matneg(a, p)}}\quad\textit{functoid}\\
= matof(nth(1, size(p)), nth(2, size(p)), vnb-lambda([i, j], cartesian(interval(1, nth(1, size(p))), interval(1, ... \\
mentioned by 7 result(s)

\smallskip

\noindent\textbf{\texttt{matof(m, n, g)}}\quad\textit{functoid}\\
= iota(p, p in mat(m, n, image(g, cartesian(interval(1, m), interval(1, n)))) and forall([i in interval(1, m), j ... \\
mentioned by 3 result(s)

\smallskip

\noindent\textbf{\texttt{matrix}}\quad\textit{operator}\\
a kernel term-former \\
mentioned by 3 result(s)

\smallskip

\noindent\textbf{\texttt{matscale(a, r, p)}}\quad\textit{functoid}\\
= matof(nth(1, size(p)), nth(2, size(p)), vnb-lambda([i, j], cartesian(interval(1, nth(1, size(p))), interval(1, ... \\
mentioned by 5 result(s)

\smallskip

\noindent\textbf{\texttt{matunit(a, n, k, l)}}\quad\textit{functoid}\\
= matof(n, n, vnb-lambda([i, j], cartesian(interval(1, n), interval(1, n)), if(i = k and j = l, one(a), zero(a)) ... \\
mentioned by 10 result(s)

\smallskip

\noindent\textbf{\texttt{metric-space}}\quad\textit{structure}\\
predicate \texttt{is-metric-space}; slots: pts dist

\smallskip

\noindent\textbf{\texttt{metric-sym-eq!}}\quad\textit{tactic}\\
\texttt{(metric-sym-eq! S P Q)} -- Add (= ((DIST S) P Q) ((DIST S) Q P)) to context via the metric-sym axiom.  [proof-local] \\
declared in \texttt{tactics-help.scm}

\smallskip

\noindent\textbf{\texttt{metric-top(md)}}\quad\textit{functoid}\\
reads: the metric topology of \$1 \\
= [pts(md), \{u in power(pts(md)): is-open(md, u)\}] \\
mentioned by 12 result(s)

\smallskip

\noindent\textbf{\texttt{metrizable-top-space}}\quad\textit{refinement}\\
predicate \texttt{is-metrizable-top-space}; same shape as top-space

\smallskip

\noindent\textbf{\texttt{minimize!}}\quad\textit{tactic}\\
\texttt{(minimize! '(v1 ... vk) GUARD MEASURE)} -- Choose v1..vk satisfying GUARD with the NN-valued MEASURE as small as possible: lands the witnesses and their minimality, and opens the two obligations any minimisation owes (MEASURE lands in NN; GUARD is satisfiable). \\
declared in \texttt{tactics-help.scm}

\smallskip

\noindent\textbf{\texttt{minor(s, r, c, n)}}\quad\textit{functoid}\\
= matof(n, n, vnb-lambda([i, j], cartesian(interval(1, n), interval(1, n)), entry(s, if(i < r, i, succ(i)), if(j ... \\
mentioned by 3 result(s)

\smallskip

\noindent\textbf{\texttt{module}}\quad\textit{structure}\\
predicate \texttt{is-module}; slots: scal vec vadd vzero vneg act

\smallskip

\noindent\textbf{\texttt{module-vector-ag(s)}}\quad\textit{view}\\
a module viewed as an abelian-group \\
mentioned by 36 result(s)

\smallskip

\noindent\textbf{\texttt{monalg(a, m)}}\quad\textit{functoid}\\
reads: the monoid algebra \$1[\$2] \\
= [finsupp(a, m), vnb-lambda([f, g], cartesian(finsupp(a, m), finsupp(a, m)), monalg-add(a, m, f, g)), vnb-lambd ... \\
mentioned by 2 result(s)

\smallskip

\noindent\textbf{\texttt{monalg-add(a, m, f, g)}}\quad\textit{functoid}\\
= vnb-lambda(x\_, carr(m), (add(a))(f(x\_), g(x\_))) \\
mentioned by 3 result(s)

\smallskip

\noindent\textbf{\texttt{monalg-mul(a, m, f, g)}}\quad\textit{functoid}\\
= vnb-lambda(x\_, carr(m), finsum(ring-additive-ag(a), vnb-lambda(p, \{p in cartesian(supp(a, m, f), supp(a, m, g) ... \\
mentioned by 7 result(s)

\smallskip

\noindent\textbf{\texttt{monalg-neg(a, m, f)}}\quad\textit{functoid}\\
= vnb-lambda(x\_, carr(m), (neg(a))(f(x\_)))

\smallskip

\noindent\textbf{\texttt{monalg-one(a, m)}}\quad\textit{functoid}\\
= vnb-lambda(x\_, carr(m), if(x\_ = iden(m), one(a), zero(a))) \\
mentioned by 2 result(s)

\smallskip

\noindent\textbf{\texttt{monalg-zero(a, m)}}\quad\textit{functoid}\\
= vnb-lambda(x\_, carr(m), zero(a))

\smallskip

\noindent\textbf{\texttt{monoid}}\quad\textit{structure}\\
predicate \texttt{is-monoid}; slots: carr opr iden

\smallskip

\noindent\textbf{\texttt{mpow}}\quad\textit{defined-fn}\\
defined by \texttt{mpow-zero} \texttt{mpow-succ} \\
mentioned by 51 result(s)

\smallskip

\noindent\textbf{\texttt{mul(s)}}\quad\textit{accessor}\\
a structure slot \\
mentioned by 239 result(s)

\smallskip


\section*{N}
\addcontentsline{toc}{section}{N}

\noindent\textbf{\texttt{nag-metric-space(nag)}}\quad\textit{functoid}\\
= [carr(nag), vnb-lambda([u, v], cartesian(carr(nag), carr(nag)), (nrm(nag))((opr(nag))(u, (inv(nag))(v))))] \\
mentioned by 11 result(s)

\smallskip

\noindent\textbf{\texttt{neg(s)}}\quad\textit{accessor}\\
a structure slot \\
mentioned by 101 result(s)

\smallskip

\noindent\textbf{\texttt{neg-rr}}\quad\textit{predicate}\\
reads: a negative real \\
a predicate

\smallskip

\noindent\textbf{\texttt{nf-metric-space(nf)}}\quad\textit{functoid}\\
= [carr(nf), vnb-lambda([x, y], cartesian(carr(nf), carr(nf)), (fnrm(nf))((add(nf))(x, (neg(nf))(y))))] \\
mentioned by 5 result(s)

\smallskip

\noindent\textbf{\texttt{ni}}\quad\textit{tactic}\\
\texttt{(ni)} -- Natural-number induction on the goal's leading FORALL over NN. \\
declared in \texttt{tactics-help.scm}

\smallskip

\noindent\textbf{\texttt{nn-add-monoid}}\quad\textit{instance}\\
an instance of comm-monoid; predicate \texttt{is-nn-add-monoid}; same shape as comm-monoid

\smallskip

\noindent\textbf{\texttt{nn-enum(s)}}\quad\textit{functoid}\\
= choice(\{f in fun(nn, s): forall([m in nn, n in nn], m < n implies f(m) < f(n))\}) \\
mentioned by 1 result(s)

\smallskip

\noindent\textbf{\texttt{nn-minus}}\quad\textit{defined-fn}\\
defined by \texttt{nn-minus-def} \\
mentioned by 11 result(s)

\smallskip

\noindent\textbf{\texttt{nnfst(n\_)}}\quad\textit{functoid}\\
reads: the first component of \$1 \\
= iota(a\_, a\_ in nn and forsome([b\_ in nn], nnpair(a\_, b\_) = n\_)) \\
mentioned by 3 result(s)

\smallskip

\noindent\textbf{\texttt{nnpair(i\_, j\_)}}\quad\textit{functoid}\\
reads: the Cantor code of (\$1, \$2) \\
= trinum(i\_ + j\_) + j\_ \\
mentioned by 7 result(s)

\smallskip

\noindent\textbf{\texttt{nnsnd(n\_)}}\quad\textit{functoid}\\
reads: the second component of \$1 \\
= iota(b\_, b\_ in nn and forsome([a\_ in nn], nnpair(a\_, b\_) = n\_)) \\
mentioned by 3 result(s)

\smallskip

\noindent\textbf{\texttt{non-zero(s)}}\quad\textit{accessor}\\
a structure slot \\
mentioned by 23 result(s)

\smallskip

\noindent\textbf{\texttt{nonneg-rr}}\quad\textit{predicate}\\
reads: a nonnegative real \\
a predicate

\smallskip

\noindent\textbf{\texttt{normed-ag}}\quad\textit{structure}\\
predicate \texttt{is-normed-ag}; slots: carr opr iden inv nrm

\smallskip

\noindent\textbf{\texttt{normed-ag-as-abelian-group(s)}}\quad\textit{view}\\
a normed-ag viewed as an abelian-group \\
mentioned by 7 result(s)

\smallskip

\noindent\textbf{\texttt{normed-field}}\quad\textit{structure}\\
predicate \texttt{is-normed-field}; slots: carr add mul neg zero one fnrm

\smallskip

\noindent\textbf{\texttt{normed-field-additive-ag(s)}}\quad\textit{view}\\
a normed-field viewed as an abelian-group \\
mentioned by 3 result(s)

\smallskip

\noindent\textbf{\texttt{normed-field-as-commutative-ring(s)}}\quad\textit{view}\\
a normed-field viewed as a commutative-ring \\
mentioned by 6 result(s)

\smallskip

\noindent\textbf{\texttt{normed-field-as-integral-domain(s)}}\quad\textit{view}\\
a normed-field viewed as an integral-domain \\
mentioned by 2 result(s)

\smallskip

\noindent\textbf{\texttt{normed-vector-space}}\quad\textit{structure}\\
predicate \texttt{is-normed-vector-space}; slots: scal vec vadd vzero vneg act vnrm

\smallskip

\noindent\textbf{\texttt{normed-vector-space-as-module(s)}}\quad\textit{view}\\
a normed-vector-space viewed as a module \\
mentioned by 4 result(s)

\smallskip

\noindent\textbf{\texttt{normed-vector-space-as-normed-ag(s)}}\quad\textit{view}\\
a normed-vector-space viewed as a normed-ag \\
mentioned by 9 result(s)

\smallskip

\noindent\textbf{\texttt{npe(m, s, f, t, g)}}\quad\textit{predicate}\\
reads: \$5 is a norm-preserving extension of \$3 from \$2 to \$4, in \$1 \\
defined by \texttt{npe} \\
mentioned by 4 result(s)

\smallskip

\noindent\textbf{\texttt{nrm(s)}}\quad\textit{accessor}\\
a structure slot \\
mentioned by 15 result(s)

\smallskip

\noindent\textbf{\texttt{nth}}\quad\textit{operator}\\
a kernel term-former \\
mentioned by 21 result(s)

\smallskip

\noindent\textbf{\texttt{nth-deriv}}\quad\textit{defined-fn}\\
defined by \texttt{nth-deriv-zero} \texttt{nth-deriv-succ} \\
mentioned by 15 result(s)

\smallskip

\noindent\textbf{\texttt{nth-deriv-v}}\quad\textit{defined-fn}\\
defined by \texttt{nth-deriv-v-zero} \texttt{nth-deriv-v-succ} \\
mentioned by 15 result(s)

\smallskip

\noindent\textbf{\texttt{nth-r}}\quad\textit{tactic}\\
\texttt{(nth-r)} -- Reduce an NTH applied to a literal LIST in the goal. \\
declared in \texttt{tactics-help.scm}

\smallskip

\noindent\textbf{\texttt{null-rr-seq(rad)}}\quad\textit{predicate}\\
reads: \$1 is a sequence of positive reals tending to zero \\
defined by \texttt{null-rr-seq} \\
mentioned by 5 result(s)

\smallskip

\noindent\textbf{\texttt{nvs-metric-space(m)}}\quad\textit{functoid}\\
= [vec(m), vnb-lambda([x, y], cartesian(vec(m), vec(m)), (vnrm(m))((vadd(m))(x, (vneg(m))(y))))] \\
mentioned by 5 result(s)

\smallskip


\section*{O}
\addcontentsline{toc}{section}{O}

\noindent\textbf{\texttt{obtain}}\quad\textit{tactic}\\
\texttt{(obtain LANE)} -- Run LANE -- a thunk whose forward step lands a `there exists' into context -- then eliminate that existential and RETURN the fresh witness's name, read off the proof state rather than guessed. \\
declared in \texttt{tactics-help.scm}

\smallskip

\noindent\textbf{\texttt{oi-l}}\quad\textit{tactic}\\
\texttt{(oi-l)} -- OR-intro left: reduce an (OR a b) goal to a. \\
declared in \texttt{tactics-help.scm}

\smallskip

\noindent\textbf{\texttt{oi-r}}\quad\textit{tactic}\\
\texttt{(oi-r)} -- OR-intro right: reduce an (OR a b) goal to b. \\
declared in \texttt{tactics-help.scm}

\smallskip

\noindent\textbf{\texttt{one(s)}}\quad\textit{accessor}\\
a structure slot \\
mentioned by 153 result(s)

\smallskip

\noindent\textbf{\texttt{opens(s)}}\quad\textit{accessor}\\
reads: the topology of \$1 \\
a structure slot \\
mentioned by 8 result(s)

\smallskip

\noindent\textbf{\texttt{opr(s)}}\quad\textit{accessor}\\
a structure slot \\
mentioned by 157 result(s)

\smallskip

\noindent\textbf{\texttt{ord-segment}}\quad\textit{operator}\\
a kernel term-former \\
mentioned by 49 result(s)

\smallskip

\noindent\textbf{\texttt{orelse}}\quad\textit{tactic}\\
\texttt{(orelse t1 t2 ...)} -- Run thunks in order, stop at the first that makes progress (LCF ORELSE). \\
declared in \texttt{tactics-help.scm}

\smallskip


\section*{P}
\addcontentsline{toc}{section}{P}

\noindent\textbf{\texttt{pair}}\quad\textit{operator}\\
a kernel term-former \\
mentioned by 38 result(s)

\smallskip

\noindent\textbf{\texttt{partial-fun}}\quad\textit{operator}\\
a kernel term-former \\
mentioned by 5 result(s)

\smallskip

\noindent\textbf{\texttt{pbc}}\quad\textit{tactic}\\
\texttt{(pbc)} -- Proof by contradiction: assume the goal's negation, prove FALSITY. \\
declared in \texttt{tactics-help.scm}

\smallskip

\noindent\textbf{\texttt{permutations(n)}}\quad\textit{functoid}\\
= injection(ord-segment(n), ord-segment(n)) \\
mentioned by 4 result(s)

\smallskip

\noindent\textbf{\texttt{pid}}\quad\textit{refinement}\\
predicate \texttt{is-pid}; same shape as integral-domain

\smallskip

\noindent\textbf{\texttt{pointwise-bounded(s, fam)}}\quad\textit{predicate}\\
reads: \$2 is pointwise bounded on \$1 \\
defined by \texttt{pointwise-bounded} \\
mentioned by 3 result(s)

\smallskip

\noindent\textbf{\texttt{poly(a)}}\quad\textit{functoid}\\
reads: \$1[x] \\
= monalg(a, nn-add-monoid) \\
mentioned by 1 result(s)

\smallskip

\noindent\textbf{\texttt{pos-rr(r)}}\quad\textit{predicate}\\
reads: a positive real \\
defined by \texttt{pos-rr} \\
mentioned by 50 result(s)

\smallskip

\noindent\textbf{\texttt{power}}\quad\textit{operator}\\
a kernel term-former \\
mentioned by 75 result(s)

\smallskip

\noindent\textbf{\texttt{preimage(s, f, v)}}\quad\textit{functoid}\\
= \{a in pts(s): f(a) in v\} \\
mentioned by 14 result(s)

\smallskip

\noindent\textbf{\texttt{prep}}\quad\textit{tactic}\\
\texttt{(prep 'ineq)} -- Why will a tactic not fire here, and what must be done first?  Runs the tactic's OWN preconditions one at a time, marks each ok / PREP / STOP, names the library lemma that repairs each unmet one, and prints a PLAN it has checked by running it on a scratch clone.  READ-ONLY: the live proof is never touched.  (prep 'ineq) is the only method so far. \\
declared in \texttt{tactics-help.scm}

\smallskip

\noindent\textbf{\texttt{principal-ideal(s, a)}}\quad\textit{functoid}\\
= \{x in carr(s): forsome([r in carr(s)], x = (mul(s))(r, a))\} \\
mentioned by 6 result(s)

\smallskip

\noindent\textbf{\texttt{prod-ord}}\quad\textit{defined-fn}\\
defined by \texttt{prod-ord-zero} \texttt{prod-ord-succ} \\
mentioned by 7 result(s)

\smallskip

\noindent\textbf{\texttt{prod-ring(r, f, s)}}\quad\textit{functoid}\\
= finprod(commutative-ring-multiplicative-cm(r), f, s) \\
mentioned by 9 result(s)

\smallskip

\noindent\textbf{\texttt{prod-set}}\quad\textit{operator}\\
a kernel term-former \\
mentioned by 7 result(s)

\smallskip

\noindent\textbf{\texttt{product-carrier(ms)}}\quad\textit{functoid}\\
= \{x in fun(nn, big-union(n, nn, pts(ms(n)))): forall([n in nn], x(n) in pts(ms(n)))\} \\
mentioned by 7 result(s)

\smallskip

\noindent\textbf{\texttt{product-metric(ms)}}\quad\textit{functoid}\\
= product-metric-w(ms, vnb-lambda(n, nn, /(1, 2 \textasciicircum{} (n + 1)))) \\
mentioned by 2 result(s)

\smallskip

\noindent\textbf{\texttt{product-metric-w(ms, w)}}\quad\textit{functoid}\\
= [product-carrier(ms), vnb-lambda([x, y], cartesian(product-carrier(ms), product-carrier(ms)), iota(l, series-c ... \\
mentioned by 7 result(s)

\smallskip

\noindent\textbf{\texttt{product-proj(ms, n)}}\quad\textit{functoid}\\
= vnb-lambda(x, product-carrier(ms), x(n)) \\
mentioned by 1 result(s)

\smallskip

\noindent\textbf{\texttt{proj(s)}}\quad\textit{functoid}\\
= vnb-lambda(a, pts(s), class(s, a)) \\
mentioned by 1 result(s)

\smallskip

\noindent\textbf{\texttt{ps-absolutely-converges-at(coef, x)}}\quad\textit{predicate}\\
reads: the power series with coefficients \$1 converges absolutely at \$2 \\
defined by \texttt{ps-absolutely-converges-at} \\
mentioned by 3 result(s)

\smallskip

\noindent\textbf{\texttt{ps-converges-at(coef, x)}}\quad\textit{predicate}\\
reads: the power series with coefficients \$1 converges at \$2 \\
defined by \texttt{ps-converges-at} \\
mentioned by 8 result(s)

\smallskip

\noindent\textbf{\texttt{ps-converges-to-at(coef, x, l)}}\quad\textit{predicate}\\
reads: the power series with coefficients \$1 converges to \$3 at \$2 \\
defined by \texttt{ps-converges-to-at} \\
mentioned by 2 result(s)

\smallskip

\noindent\textbf{\texttt{ps-partial-sum(coef, x, k)}}\quad\textit{functoid}\\
= sum-ag(normed-field-additive-ag(rr-normed-field), vnb-lambda(n, nn, coef(n) * x \textasciicircum{} n), k) \\
mentioned by 6 result(s)

\smallskip

\noindent\textbf{\texttt{ps-ratio-limit(coef, l)}}\quad\textit{predicate}\\
reads: the coefficients \$1 have ratio limit \$2 \\
defined by \texttt{ps-ratio-limit} \\
mentioned by 3 result(s)

\smallskip

\noindent\textbf{\texttt{pseudo-gauge-top(fam, ground)}}\quad\textit{functoid}\\
reads: the topology generated by the pseudometric family \$1 on \$2 \\
= [ground, \{u in power(ground): is-open-gauge(fam, ground, u)\}] \\
mentioned by 2 result(s)

\smallskip

\noindent\textbf{\texttt{pseudometric-space}}\quad\textit{structure}\\
predicate \texttt{is-pseudometric-space}; slots: pts dist

\smallskip

\noindent\textbf{\texttt{pts(s)}}\quad\textit{accessor}\\
a structure slot \\
mentioned by 165 result(s)

\smallskip


\section*{Q}
\addcontentsline{toc}{section}{Q}

\noindent\textbf{\texttt{qed}}\quad\textit{tactic}\\
\texttt{(qed 'name)} -- Install the finished proof as theorem `name' and save its replayable script. \\
declared in \texttt{tactics-help.scm}

\smallskip

\noindent\textbf{\texttt{qq-field}}\quad\textit{instance}\\
an instance of field; predicate \texttt{is-qq-field}; same shape as field

\smallskip

\noindent\textbf{\texttt{qq-ring}}\quad\textit{instance}\\
an instance of ring; predicate \texttt{is-qq-ring}; same shape as integral-domain

\smallskip

\noindent\textbf{\texttt{qrfl}}\quad\textit{tactic}\\
\texttt{(qrfl)} -- Quasi-reflexivity: close t = t under the partial-equality definedness reading. \\
declared in \texttt{tactics-help.scm}

\smallskip

\noindent\textbf{\texttt{quietly}}\quad\textit{tactic}\\
\texttt{(quietly thunk)} -- Run thunk with state-dump output suppressed; returns its value. \\
declared in \texttt{tactics-help.scm}

\smallskip

\noindent\textbf{\texttt{quotient(s)}}\quad\textit{functoid}\\
= image(proj(s), pts(s)) \\
mentioned by 9 result(s)

\smallskip


\section*{R}
\addcontentsline{toc}{section}{R}

\noindent\textbf{\texttt{ran(f)}}\quad\textit{functoid}\\
= image(f, dom(f)) \\
mentioned by 7 result(s)

\smallskip

\noindent\textbf{\texttt{real-part}}\quad\textit{operator}\\
a kernel term-former

\smallskip

\noindent\textbf{\texttt{recip(s)}}\quad\textit{accessor}\\
a structure slot \\
mentioned by 45 result(s)

\smallskip

\noindent\textbf{\texttt{reduce}}\quad\textit{defined-fn}\\
defined by \texttt{reduce-one} \texttt{reduce-succ} \\
mentioned by 26 result(s)

\smallskip

\noindent\textbf{\texttt{rel(s)}}\quad\textit{accessor}\\
a structure slot \\
mentioned by 4 result(s)

\smallskip

\noindent\textbf{\texttt{rel-free(md, n, u)}}\quad\textit{predicate}\\
reads: the \$2 vectors \$3 are linearly independent in \$1 \\
defined by \texttt{rel-free} \\
mentioned by 5 result(s)

\smallskip

\noindent\textbf{\texttt{related(s, a, b)}}\quad\textit{functoid}\\
= [a, b] in rel(s) \\
mentioned by 2 result(s)

\smallskip

\noindent\textbf{\texttt{repeat}}\quad\textit{tactic}\\
\texttt{(repeat thunk [cap])} -- Run thunk until it stops changing the proof state (LCF REPEAT). \\
declared in \texttt{tactics-help.scm}

\smallskip

\noindent\textbf{\texttt{replay-proof}}\quad\textit{tactic}\\
\texttt{(replay-proof 'name [subst])} -- Re-run a saved script on the current goal, optionally renaming free vars. \\
declared in \texttt{tactics-help.scm}

\smallskip

\noindent\textbf{\texttt{res}}\quad\textit{operator}\\
a kernel term-former \\
mentioned by 4 result(s)

\smallskip

\noindent\textbf{\texttt{respects(s, f)}}\quad\textit{functoid}\\
= forall([a in pts(s), b in pts(s)], related(s, a, b) implies f(a) = f(b)) \\
mentioned by 4 result(s)

\smallskip

\noindent\textbf{\texttt{respects2(s, f)}}\quad\textit{functoid}\\
= forall([a, b, a\_, b\_], a in pts(s) implies b in pts(s) implies a\_ in pts(s) implies b\_ in pts(s) implies relat ... \\
mentioned by 3 result(s)

\smallskip

\noindent\textbf{\texttt{restvar}}\quad\textit{operator}\\
a kernel term-former \\
mentioned by 4 result(s)

\smallskip

\noindent\textbf{\texttt{rfl}}\quad\textit{tactic}\\
\texttt{(rfl)} -- Close a reflexive equality goal (t = t). \\
declared in \texttt{tactics-help.scm}

\smallskip

\noindent\textbf{\texttt{ring}}\quad\textit{structure}\\
predicate \texttt{is-ring}; slots: carr add mul neg zero one

\smallskip

\noindent\textbf{\texttt{ring-additive-ag(s)}}\quad\textit{view}\\
a ring viewed as an abelian-group \\
mentioned by 151 result(s)

\smallskip

\noindent\textbf{\texttt{ring-multiplicative-monoid(s)}}\quad\textit{view}\\
a ring viewed as a monoid \\
mentioned by 4 result(s)

\smallskip

\noindent\textbf{\texttt{ring-power(r, x, n)}}\quad\textit{functoid}\\
= mpow(commutative-ring-multiplicative-cm(r), x, n) \\
mentioned by 13 result(s)

\smallskip

\noindent\textbf{\texttt{ring-prod(x, y)}}\quad\textit{functoid}\\
= [cartesian(carr(x), carr(y)), vnb-lambda([p, q], cartesian(cartesian(carr(x), carr(y)), cartesian(carr(x), car ... \\
mentioned by 4 result(s)

\smallskip

\noindent\textbf{\texttt{ring-prod-n}}\quad\textit{defined-fn}\\
defined by \texttt{ring-prod-n-zero} \texttt{ring-prod-n-succ} \\
mentioned by 5 result(s)

\smallskip

\noindent\textbf{\texttt{ringoid}}\quad\textit{structure}\\
predicate \texttt{is-ringoid}; slots: carr add mul neg zero one idl

\smallskip

\noindent\textbf{\texttt{ringoid-as-ring(s)}}\quad\textit{view}\\
a ringoid viewed as a ring \\
mentioned by 4 result(s)

\smallskip

\noindent\textbf{\texttt{ringoid-quotient(r)}}\quad\textit{functoid}\\
= quotient(ringoid-setoid(r))

\smallskip

\noindent\textbf{\texttt{ringoid-quotient-ring(r)}}\quad\textit{functoid}\\
= [quotient(ringoid-setoid(r)), descend2(ringoid-setoid(r), vnb-lambda([a, b], cartesian(carr(r), carr(r)), clas ... \\
mentioned by 15 result(s)

\smallskip

\noindent\textbf{\texttt{ringoid-rel(r)}}\quad\textit{functoid}\\
= \{p in cartesian(carr(r), carr(r)): (add(r))(nth(1, p), (neg(r))(nth(2, p))) in idl(r)\} \\
mentioned by 3 result(s)

\smallskip

\noindent\textbf{\texttt{ringoid-setoid(r)}}\quad\textit{functoid}\\
= [carr(r), ringoid-rel(r)] \\
mentioned by 17 result(s)

\smallskip

\noindent\textbf{\texttt{rpow}}\quad\textit{operator}\\
a kernel term-former \\
mentioned by 27 result(s)

\smallskip

\noindent\textbf{\texttt{rr+*-add-monoid}}\quad\textit{refinement}\\
predicate \texttt{is-rr+*-add-monoid}; same shape as comm-monoid

\smallskip

\noindent\textbf{\texttt{rr-bounded-above(s)}}\quad\textit{predicate}\\
reads: \$1 is bounded above \\
defined by \texttt{rr-bounded-above} \\
mentioned by 5 result(s)

\smallskip

\noindent\textbf{\texttt{rr-bounded-ms(())}}\quad\textit{functoid}\\
= bdd-metric(rr-ms)

\smallskip

\noindent\textbf{\texttt{rr-ms}}\quad\textit{instance}\\
an instance of metric-space; predicate \texttt{is-rr-ms}; same shape as metric-space

\smallskip

\noindent\textbf{\texttt{rr-normed-field}}\quad\textit{instance}\\
an instance of normed-field; predicate \texttt{is-rr-normed-field}; same shape as normed-field

\smallskip

\noindent\textbf{\texttt{rr-upper-bound(s, b)}}\quad\textit{predicate}\\
reads: \$2 is an upper bound of \$1 \\
defined by \texttt{rr-upper-bound} \\
mentioned by 6 result(s)

\smallskip

\noindent\textbf{\texttt{rs}}\quad\textit{tactic}\\
\texttt{(rs)} -- Ring-simplify the goal (normal form over the ambient ring). \\
declared in \texttt{tactics-help.scm}

\smallskip


\section*{S}
\addcontentsline{toc}{section}{S}

\noindent\textbf{\texttt{save-proof}}\quad\textit{tactic}\\
\texttt{(save-proof 'name)} -- Save the current script under a name without finishing. \\
declared in \texttt{tactics-help.scm}

\smallskip

\noindent\textbf{\texttt{scal(s)}}\quad\textit{accessor}\\
a structure slot \\
mentioned by 125 result(s)

\smallskip

\noindent\textbf{\texttt{scout}}\quad\textit{tactic}\\
\texttt{(scout [depth [branch [nodes]]])} -- Speculatively search to a bounded depth across INDEPENDENT scratch branches; RETURNS a nested list (number-of-branches (d b) goal best-partials closing-branches) rather than printing.  Never touches the live proof. \\
declared in \texttt{tactics-help.scm}

\smallskip

\noindent\textbf{\texttt{scout-run}}\quad\textit{tactic}\\
\texttt{(scout-run k)} -- Adopt closing branch k from the last (scout)/(scout-show) onto the live proof, running its tactics for real (they record and display normally). \\
declared in \texttt{tactics-help.scm}

\smallskip

\noindent\textbf{\texttt{scout-show}}\quad\textit{tactic}\\
\texttt{(scout-show [depth [branch [nodes]]])} -- Like (scout), but PRINTS the human-readable report (examined count, goal, closing branches or best partials) to the REPL.  Returns the same nested list (scout) does. \\
declared in \texttt{tactics-help.scm}

\smallskip

\noindent\textbf{\texttt{semigroup}}\quad\textit{structure}\\
predicate \texttt{is-semigroup}; slots: carr opr

\smallskip

\noindent\textbf{\texttt{sep}}\quad\textit{operator}\\
a kernel term-former \\
mentioned by 9 result(s)

\smallskip

\noindent\textbf{\texttt{sep-me}}\quad\textit{tactic}\\
\texttt{(sep-me hyp)} -- Separation membership elim: split a separation-membership assumption into A-membership and the predicate. \\
declared in \texttt{tactics-help.scm}

\smallskip

\noindent\textbf{\texttt{sep-mi}}\quad\textit{tactic}\\
\texttt{(sep-mi)} -- Separation membership intro: prove (IN t \{x in A | p\}). \\
declared in \texttt{tactics-help.scm}

\smallskip

\noindent\textbf{\texttt{sep-set}}\quad\textit{tactic}\\
\texttt{(sep-set)} -- Separation sethood: the separation set \{x in A | p\} is a set. \\
declared in \texttt{tactics-help.scm}

\smallskip

\noindent\textbf{\texttt{seq-compact(s)}}\quad\textit{predicate}\\
reads: \$1 is sequentially compact \\
defined by \texttt{seq-compact} \\
mentioned by 7 result(s)

\smallskip

\noindent\textbf{\texttt{series-converges(f)}}\quad\textit{predicate}\\
reads: the series \$1 converges \\
defined by \texttt{series-converges} \\
mentioned by 11 result(s)

\smallskip

\noindent\textbf{\texttt{series-converges-to(f, l)}}\quad\textit{predicate}\\
reads: the series \$1 converges to \$2 \\
defined by \texttt{series-converges-to} \\
mentioned by 3 result(s)

\smallskip

\noindent\textbf{\texttt{series-partial-sum(f, k)}}\quad\textit{functoid}\\
= sum-ag(normed-field-additive-ag(rr-normed-field), f, k) \\
mentioned by 12 result(s)

\smallskip

\noindent\textbf{\texttt{setoid}}\quad\textit{structure}\\
predicate \texttt{is-setoid}; slots: pts rel

\smallskip

\noindent\textbf{\texttt{show}}\quad\textit{tactic}\\
\texttt{(show)} -- Redisplay the current proof state. \\
declared in \texttt{tactics-help.scm}

\smallskip

\noindent\textbf{\texttt{simp}}\quad\textit{tactic}\\
\texttt{(simp [target])} -- Rewrite a commutative-ring SUBTERM of the goal to canonical form, IN PLACE (e.g. (x+y)\textasciicircum{}2 inside a larger goal becomes x\textasciicircum{}2 + 2*x*y + y\textasciicircum{}2).  Works on BOTH surfaces: concrete number domains (+ * - \textasciicircum{} over NN/ZZ/QQ/RR/CC) and a generic ring s ((ADD s)/(MUL s)/(NEG s), carrier (CARR s)).  No arg = outermost ring subterm; "term" targets a specific one.  Sound by cut + crs + eq-subst (no new kernel rule); needs the subterm's generators typed in context (true post-di), else refuses and names them. \\
declared in \texttt{tactics-help.scm}

\smallskip

\noindent\textbf{\texttt{sin}}\quad\textit{operator}\\
a kernel term-former

\smallskip

\noindent\textbf{\texttt{singleton}}\quad\textit{operator}\\
a kernel term-former \\
mentioned by 8 result(s)

\smallskip

\noindent\textbf{\texttt{size(m)}}\quad\textit{functoid}\\
= [length(m), length(nth(1, m))]

\smallskip

\noindent\textbf{\texttt{slot}}\quad\textit{tactic}\\
\texttt{(slot 'acc)} -- Reduce a structure ACCESSOR to its projection: (CARR s) becomes (NTH 1 s).  The one door for accessor reductions -- use it, never mac, on an accessor name. \\
declared in \texttt{tactics-help.scm}

\smallskip

\noindent\textbf{\texttt{smith-staircase(a, m, n, d, k)}}\quad\textit{predicate}\\
reads: \$4 is in Smith staircase form of rank \$5, as an \$2-by-\$3 matrix over \$1 \\
defined by \texttt{smith-staircase} \\
mentioned by 3 result(s)

\smallskip

\noindent\textbf{\texttt{snoc-col(w, n, x)}}\quad\textit{functoid}\\
= matof(succ(n), 1, vnb-lambda([i\_, j\_], cartesian(interval(1, succ(n)), interval(1, 1)), if(i\_ = succ(n), x, en ... \\
mentioned by 7 result(s)

\smallskip

\noindent\textbf{\texttt{snoc-row(c, n, r)}}\quad\textit{functoid}\\
= matof(1, succ(n), vnb-lambda([i\_, j\_], cartesian(interval(1, 1), interval(1, succ(n))), if(j\_ = succ(n), r, en ... \\
mentioned by 7 result(s)

\smallskip

\noindent\textbf{\texttt{sos}}\quad\textit{tactic}\\
\texttt{(sos "c1" "c2" ...)} -- Sum-of-squares closer for a nonstrict polynomial inequality a <= b over RR (the nonlinear companion of (ineq)).  You supply the terms to be SQUARED; it finds the nonnegative coefficients.  (tactics 'sos) for the worked example. \\
declared in \texttt{tactics-help.scm}

\smallskip

\noindent\textbf{\texttt{sp}}\quad\textit{tactic}\\
\texttt{(sp goal)} -- Start a proof of `goal' (a "string", a raw S-expr, or a wff); clears the script. \\
declared in \texttt{tactics-help.scm}

\smallskip

\noindent\textbf{\texttt{span(md, n, u)}}\quad\textit{functoid}\\
= \{x\_ in vec(md): forsome([c\_ in mat(1, n, carr(scal(md)))], x\_ = entry(matact(md, c\_, u), 1, 1))\} \\
mentioned by 6 result(s)

\smallskip

\noindent\textbf{\texttt{span-add-one(m, s, v)}}\quad\textit{functoid}\\
= \{y\_ in vec(m): forsome([x\_ in s, r\_ in rr], y\_ = (vadd(m))(x\_, (act(m))(r\_, v)))\} \\
mentioned by 8 result(s)

\smallskip

\noindent\textbf{\texttt{spans(md, n, u, sm)}}\quad\textit{predicate}\\
reads: the \$2 vectors \$3 span \$4 in \$1 \\
defined by \texttt{spans} \\
mentioned by 10 result(s)

\smallskip

\noindent\textbf{\texttt{splice}}\quad\textit{operator}\\
a kernel term-former \\
mentioned by 4 result(s)

\smallskip

\noindent\textbf{\texttt{split-ands!}}\quad\textit{tactic}\\
\texttt{(split-ands!)} -- Flatten every AND assumption of the focus into separate assumptions.  [proof-local] \\
declared in \texttt{tactics-help.scm}

\smallskip

\noindent\textbf{\texttt{sqrt}}\quad\textit{operator}\\
a kernel term-former \\
mentioned by 13 result(s)

\smallskip

\noindent\textbf{\texttt{strictly-mono-nn(phi)}}\quad\textit{predicate}\\
reads: \$1 is a strictly increasing sequence of naturals \\
defined by \texttt{strictly-mono-nn} \\
mentioned by 17 result(s)

\smallskip

\noindent\textbf{\texttt{submat(s, p, q)}}\quad\textit{functoid}\\
= matof(p, q, vnb-lambda([i, j], cartesian(interval(1, p), interval(1, q)), entry(s, succ(i), succ(j)))) \\
mentioned by 6 result(s)

\smallskip

\noindent\textbf{\texttt{subseq(f, phi)}}\quad\textit{functoid}\\
= vnb-lambda(k, nn, f(phi(k))) \\
mentioned by 12 result(s)

\smallskip

\noindent\textbf{\texttt{subset}}\quad\textit{primitive}\\
reads: \$1 is a subset of \$2 \\
a kernel relation \\
mentioned by 97 result(s)

\smallskip

\noindent\textbf{\texttt{subst}}\quad\textit{tactic}\\
\texttt{(subst '(= s t))  |  (subst '(== s t))} -- Rewrite s -> t throughout the goal, using an equation s = t -- or a quasi-equation s == t -- that is in context, in EITHER orientation (Leibniz substitution).  Cannot reach a term in OPERATOR position: use the equation as a macete instead. \\
declared in \texttt{tactics-help.scm}

\smallskip

\noindent\textbf{\texttt{succ}}\quad\textit{operator}\\
a kernel term-former \\
mentioned by 228 result(s)

\smallskip

\noindent\textbf{\texttt{succ\_ord}}\quad\textit{operator}\\
a kernel term-former \\
mentioned by 23 result(s)

\smallskip

\noindent\textbf{\texttt{sum}}\quad\textit{defined-fn}\\
defined by \texttt{sum-zero} \texttt{sum-succ} \\
mentioned by 14 result(s)

\smallskip

\noindent\textbf{\texttt{sum-ag}}\quad\textit{defined-fn}\\
defined by \texttt{sum-ag-zero} \texttt{sum-ag-succ} \\
mentioned by 26 result(s)

\smallskip

\noindent\textbf{\texttt{sum-set}}\quad\textit{operator}\\
a kernel term-former \\
mentioned by 11 result(s)

\smallskip

\noindent\textbf{\texttt{summable-weight(w)}}\quad\textit{predicate}\\
reads: \$1 is a summable sequence of positive weights \\
defined by \texttt{summable-weight} \\
mentioned by 9 result(s)

\smallskip

\noindent\textbf{\texttt{sums-to(grp, f, r)}}\quad\textit{predicate}\\
reads: \$2 sums to \$3 in \$1 \\
defined by \texttt{sums-to} \\
mentioned by 6 result(s)

\smallskip

\noindent\textbf{\texttt{sup}}\quad\textit{functoid}\\
reads: the least upper bound of \$1 \\
mentioned by 3 result(s)

\smallskip

\noindent\textbf{\texttt{sup-ord}}\quad\textit{operator}\\
a kernel term-former \\
mentioned by 9 result(s)

\smallskip

\noindent\textbf{\texttt{supp(a, m, f)}}\quad\textit{functoid}\\
reads: the support of \$3 \\
= \{x\_ in carr(m): not(f(x\_) = zero(a))\}

\smallskip


\section*{T}
\addcontentsline{toc}{section}{T}

\noindent\textbf{\texttt{ta}}\quad\textit{tactic}\\
\texttt{(ta 'name)} -- Theorem-assumption: bring the named installed theorem into context as an assumption. \\
declared in \texttt{tactics-help.scm}

\smallskip

\noindent\textbf{\texttt{taylor-differentiable(f, a, x, n)}}\quad\textit{predicate}\\
reads: \$1 is \$4 times differentiable from \$2 to \$3 \\
defined by \texttt{taylor-differentiable} \\
mentioned by 9 result(s)

\smallskip

\noindent\textbf{\texttt{taylor-differentiable-v(m, f, a, x, n)}}\quad\textit{predicate}\\
reads: \$2 is \$5 times differentiable from \$3 to \$4, as a curve in \$1 \\
defined by \texttt{taylor-differentiable-v} \\
mentioned by 5 result(s)

\smallskip

\noindent\textbf{\texttt{taylor-poly(f, a, n, x)}}\quad\textit{functoid}\\
= series-partial-sum(vnb-lambda(k, nn, (nth-deriv(f, k))(a) * (x - a) \textasciicircum{} k * recip(factorial(k))), succ(n)) \\
mentioned by 15 result(s)

\smallskip

\noindent\textbf{\texttt{taylor-poly-v}}\quad\textit{defined-fn}\\
defined by \texttt{taylor-poly-v-zero} \texttt{taylor-poly-v-succ} \\
mentioned by 10 result(s)

\smallskip

\noindent\textbf{\texttt{te}}\quad\textit{tactic}\\
\texttt{(te hyp k)} -- Tuple elim: project the k-th component out of a tuple-membership assumption. \\
declared in \texttt{tactics-help.scm}

\smallskip

\noindent\textbf{\texttt{ti}}\quad\textit{tactic}\\
\texttt{(ti)} -- Tuple intro: prove a tuple/LIST membership component-wise. \\
declared in \texttt{tactics-help.scm}

\smallskip

\noindent\textbf{\texttt{top-space}}\quad\textit{structure}\\
predicate \texttt{is-top-space}; slots: pts opens

\smallskip

\noindent\textbf{\texttt{totally-bounded(s)}}\quad\textit{predicate}\\
reads: \$1 is totally bounded \\
defined by \texttt{totally-bounded} \\
mentioned by 11 result(s)

\smallskip

\noindent\textbf{\texttt{trinum}}\quad\textit{functoid}\\
reads: the \$1-th triangular number \\
defined by \texttt{trinum-zero} \texttt{trinum-succ} \\
mentioned by 8 result(s)

\smallskip

\noindent\textbf{\texttt{tuples}}\quad\textit{operator}\\
a kernel term-former \\
mentioned by 7 result(s)

\smallskip


\section*{U}
\addcontentsline{toc}{section}{U}

\noindent\textbf{\texttt{ue}}\quad\textit{tactic}\\
\texttt{(ue hyp)} -- Union elim: split a cited (IN x (UNION ...)) membership assumption into one subgoal per set -- union-side case analysis, the dual of `ui'. \\
declared in \texttt{tactics-help.scm}

\smallskip

\noindent\textbf{\texttt{ui}}\quad\textit{tactic}\\
\texttt{(ui k)} -- Union intro: reduce a goal (IN x (UNION ...)) to membership of x in the k-th set (1-based). \\
declared in \texttt{tactics-help.scm}

\smallskip

\noindent\textbf{\texttt{union}}\quad\textit{operator}\\
a kernel term-former \\
mentioned by 39 result(s)

\smallskip

\noindent\textbf{\texttt{unitrow(a, n, i)}}\quad\textit{functoid}\\
= matof(1, n, vnb-lambda([rw, cl], cartesian(interval(1, 1), interval(1, n)), if(cl = i, one(a), zero(a)))) \\
mentioned by 7 result(s)

\smallskip


\section*{V}
\addcontentsline{toc}{section}{V}

\noindent\textbf{\texttt{vadd(s)}}\quad\textit{accessor}\\
a structure slot \\
mentioned by 268 result(s)

\smallskip

\noindent\textbf{\texttt{vec(s)}}\quad\textit{accessor}\\
a structure slot \\
mentioned by 389 result(s)

\smallskip

\noindent\textbf{\texttt{vector-space}}\quad\textit{refinement}\\
predicate \texttt{is-vector-space}; same shape as module

\smallskip

\noindent\textbf{\texttt{vlet}}\quad\textit{tactic}\\
\texttt{(vlet (n1 ...) FORMER)} -- Bind proof-object names from the current state.  FORMER is (match PATTERN) -- bind each name to its slot in a context formula matching PATTERN -- or (choice [v body]) -- eliminate an in-context existential and bind its witness. \\
declared in \texttt{tactics-help.scm}

\smallskip

\noindent\textbf{\texttt{vnb-lambda}}\quad\textit{operator}\\
a kernel term-former \\
mentioned by 217 result(s)

\smallskip

\noindent\textbf{\texttt{vneg(s)}}\quad\textit{accessor}\\
a structure slot \\
mentioned by 50 result(s)

\smallskip

\noindent\textbf{\texttt{vnrm(s)}}\quad\textit{accessor}\\
a structure slot \\
mentioned by 35 result(s)

\smallskip

\noindent\textbf{\texttt{vzero(s)}}\quad\textit{accessor}\\
a structure slot \\
mentioned by 172 result(s)

\smallskip


\section*{W}
\addcontentsline{toc}{section}{W}

\noindent\textbf{\texttt{wbc}}\quad\textit{tactic}\\
\texttt{(wbc ['name])} -- Witness-shape backchain: on an `exists v. ...' goal, cite a PSS lemma that MANUFACTURES a witness of that shape, so you can finish with inst+/grind/ew. \\
declared in \texttt{tactics-help.scm}

\smallskip

\noindent\textbf{\texttt{wk}}\quad\textit{tactic}\\
\texttt{(wk hyp)} -- Weaken: drop a cited assumption from the context to tidy the hypothesis list.  `hyp' may be a formula, a "string", or a 1-based assumption index. \\
declared in \texttt{tactics-help.scm}

\smallskip


\section*{Z}
\addcontentsline{toc}{section}{Z}

\noindent\textbf{\texttt{zero(s)}}\quad\textit{accessor}\\
a structure slot \\
mentioned by 385 result(s)

\smallskip

\noindent\textbf{\texttt{zero-ring}}\quad\textit{defined-fn}\\
defined by \texttt{zero-ring-def}

\smallskip

\noindent\textbf{\texttt{zeromat(a, m, n)}}\quad\textit{functoid}\\
= matof(m, n, vnb-lambda([i, j], cartesian(interval(1, m), interval(1, n)), zero(a))) \\
mentioned by 15 result(s)

\smallskip

\noindent\textbf{\texttt{zkept}}\quad\textit{defined-fn}\\
defined by \texttt{zkept-zero} \texttt{zkept-succ} \texttt{zkept-limit} \\
mentioned by 9 result(s)

\smallskip

\noindent\textbf{\texttt{zup}}\quad\textit{defined-fn}\\
defined by \texttt{zup-zero} \texttt{zup-succ} \texttt{zup-limit} \\
mentioned by 9 result(s)

\smallskip

\noindent\textbf{\texttt{zz-act}}\quad\textit{defined-fn}\\
defined by \texttt{zz-act-nonneg} \texttt{zz-act-neg} \\
mentioned by 52 result(s)

\smallskip

\noindent\textbf{\texttt{zz-bezout-set(a, b)}}\quad\textit{functoid}\\
= \{z in zz: forsome([x\_ in zz, y\_ in zz], z = x\_ * a + y\_ * b)\} \\
mentioned by 4 result(s)

\smallskip

\noindent\textbf{\texttt{zz-coprime(a, b)}}\quad\textit{predicate}\\
reads: \$1 and \$2 are coprime \\
defined by \texttt{zz-coprime} \\
mentioned by 2 result(s)

\smallskip

\noindent\textbf{\texttt{zz-divides(a, b)}}\quad\textit{predicate}\\
reads: \$1 divides \$2 \\
defined by \texttt{zz-divides} \\
mentioned by 5 result(s)

\smallskip

\noindent\textbf{\texttt{zz-is-gcd(a, b, d)}}\quad\textit{predicate}\\
reads: \$3 is a greatest common divisor of \$1 and \$2 \\
defined by \texttt{zz-is-gcd} \\
mentioned by 4 result(s)

\smallskip

\noindent\textbf{\texttt{zz-ring}}\quad\textit{instance}\\
an instance of ring; predicate \texttt{is-zz-ring}; same shape as euclidean-ring

\smallskip

\endgroup
